\documentclass[11pt,a4paper]{article}
\usepackage[margin=2.4cm]{geometry}
\usepackage{amsmath,amsthm,mathtools}
\usepackage{fontspec}
\usepackage{unicode-math}
\setmainfont{STIX2Text}[Path=fonts/, Extension=.otf,
  UprightFont=*-Regular, ItalicFont=*-Italic,
  BoldFont=*-Bold, BoldItalicFont=*-BoldItalic]
\setmathfont{STIX2Math}[Path=fonts/, Extension=.otf]
\usepackage[protrusion=true,expansion=false]{microtype}
\usepackage{booktabs,enumitem}
\usepackage{xcolor}
\usepackage[psdextra,hidelinks]{hyperref}
\emergencystretch=3em
\hbadness=10000
\hfuzz=2pt
\newtheorem{theorem}{Theorem}[section]
\newtheorem{lemma}[theorem]{Lemma}
\newtheorem{proposition}[theorem]{Proposition}
\newtheorem{corollary}[theorem]{Corollary}
\theoremstyle{definition}
\newtheorem{definition}[theorem]{Definition}
\newtheorem{construction}[theorem]{Construction}
\newtheorem{assumption}[theorem]{Assumption}
\newtheorem{example}[theorem]{Example}

% kind-coded theorem blocks, palette shared with the figures and the website:
% definitions/constructions blue (tierin), assumptions amber (tierfull),
% proved statements a neutral bar; remarks and examples run unframed
\usepackage{mdframed}
\providecolor{tierin}{RGB}{60,130,180}
\providecolor{tierfull}{RGB}{210,140,40}
\mdfdefinestyle{kindbar}{topline=false,bottomline=false,rightline=false,
  linewidth=1.5pt,innerleftmargin=9pt,innerrightmargin=4pt,
  innertopmargin=5pt,innerbottommargin=6pt,leftmargin=0pt,rightmargin=0pt,
  skipabove=10pt,skipbelow=10pt}
\surroundwithmdframed[style=kindbar,linecolor=tierin!70,backgroundcolor=tierin!4]{definition}
\surroundwithmdframed[style=kindbar,linecolor=tierin!70,backgroundcolor=tierin!4]{construction}
\surroundwithmdframed[style=kindbar,linecolor=tierfull!80,backgroundcolor=tierfull!5]{assumption}
\surroundwithmdframed[style=kindbar,linecolor=black!30]{theorem}
\surroundwithmdframed[style=kindbar,linecolor=black!30]{lemma}
\surroundwithmdframed[style=kindbar,linecolor=black!30]{proposition}
\surroundwithmdframed[style=kindbar,linecolor=black!30]{corollary}

\usepackage{fancyhdr}
\pagestyle{fancy}
\fancyhf{}
\fancyhead[L]{\footnotesize\scshape The Zcash Arboretum}
\fancyhead[R]{\footnotesize\itshape\nouppercase{\leftmark}}
\fancyfoot[C]{\footnotesize\thepage}
\renewcommand{\headrulewidth}{0.3pt}
\renewcommand{\sectionmark}[1]{\markboth{\thesection\ \ #1}{}}
\theoremstyle{remark}
\newtheorem{remark}[theorem]{Remark}
\newcommand{\N}{\mathbb{N}}
\newcommand{\Z}{\mathbb{Z}}
\newcommand{\Q}{\mathbb{Q}}
\newcommand{\R}{\mathbb{R}}
\newcommand{\F}{\mathbb{F}}
\newcommand{\G}{\mathbb{G}}
\newcommand{\E}{\mathbb{E}}
\newcommand{\Pallas}{\mathsf{Pallas}}
\newcommand{\Vesta}{\mathsf{Vesta}}
\newcommand{\Ep}{\mathsf{E}_{\Pallas}}
\newcommand{\Ev}{\mathsf{E}_{\Vesta}}
\newcommand{\ip}[2]{\langle #1,\, #2\rangle}
\DeclareMathOperator{\ord}{ord}
\DeclareMathOperator{\lcm}{lcm}
\DeclareMathOperator{\rank}{rank}
\DeclareMathOperator{\im}{im}
\DeclareMathOperator{\SD}{SD}
\DeclareMathOperator{\LE}{LE}
\DeclareMathOperator{\BE}{BE}
\DeclareMathOperator{\val}{val}
\DeclareMathOperator{\negl}{negl}
\usepackage[nameinlink]{cleveref}
\renewenvironment{abstract}{%
  \small\begin{center}{\bfseries\abstractname}\end{center}\medskip\noindent\ignorespaces
}{\par\medskip}

\title{\textbf{\Huge Math Guide}\\[6pt]\large Arithmetic, randomness, and algebra for shielded payments}
\author{m@rek.onl}
\date{}
\begin{document}
\maketitle
\thispagestyle{empty}
\begin{abstract}
This volume of \emph{The Zcash Arboretum} constructs the mathematical
objects used by its cryptographic volumes. It assumes elementary integer
arithmetic and the ability to follow an algorithm. We build ways to
represent data, calculate with it and reason about random choices,
introducing each required concept before applying it. Proofs establish
the stated elementary results; facts imported without proof are identified
explicitly. Knowing how a calculation works does not by itself establish
that an attacker cannot reverse it.
\end{abstract}
\clearpage
\tableofcontents
\clearpage

\section{Language, integers, and encodings}\label{sec:sets}

A calculation on numbers and a message sent between computers are different
objects. We first specify how messages represent numbers, which messages
to reject, and whether different messages can denote the same number.

\subsection{Objects, functions and assertions}

A \emph{set} is a collection whose members are distinct. We write
\(x\in S\) for membership, \(S\subseteq T\) when every member of \(S\)
belongs to \(T\), and \(\varnothing\) for the empty set. The symbol
\(S\cup T\) denotes the union: members in either set. The intersection
\(S\cap T\) contains members in both; \(S\setminus T\) contains those
in \(S\) but not \(T\). The integer \(|S|\)
counts the elements of a finite set. An ordered pair \((x,y)\) records
both its entries and their order. The product \(S\times T\) is the set
of such pairs with \(x\in S\) and \(y\in T\). The notation \(S^n\)
means ordered sequences, or tuples, of exactly \(n\) elements of \(S\),
including one empty sequence when \(n=0\).

A \emph{function} \(f:S\to T\) assigns exactly one output \(f(x)\in T\)
to each input \(x\in S\). The sets \(S,T\) are its domain and codomain,
respectively. Its image is
\(\im f=\{f(x):x\in S\}\). It is injective when equal outputs require
equal inputs, surjective when \(\im f=T\), and bijective when both
conditions hold. A bijection has an inverse function, which returns \(x\)
from \(f(x)\). The composition
\(g\circ f\) means first apply \(f\), then \(g\). A function that may
reject has codomain \(T\cup\{\bot\}\), where \(\bot\) is a distinct
failure symbol, not the number zero or an empty sequence.

An assertion has a truth value. The expressions \(A\land B\),
\(A\lor B\), and \(\neg A\) mean ``both'', ``at least one'', and
``not''. The implication \(A\Rightarrow B\) prohibits the case in which
\(A\) is true and \(B\) false. Its contrapositive
\(\neg B\Rightarrow\neg A\) has exactly the same truth value.
An equivalence \(A\Longleftrightarrow B\) asserts both implications.
The quantifiers \(\forall x\in S\) and \(\exists x\in S\) mean
``for every member'' and ``for at least one member''. Their order matters:
\(\forall x\,\exists y\) permits a different \(y\) for each \(x\).
An existence proof does not, by itself, give an efficient way to find the
object it proves exists.

The natural numbers are \(\N=\{0,1,2,\ldots\}\); the integers
\(\Z\) include their negatives. We use induction as follows. To prove
an assertion for every \(n\in\N\), prove it at zero and prove that its
truth at \(n\) implies its truth at \(n+1\). Equivalently, a nonempty
set of natural numbers has a least member: a least counterexample could
be neither zero nor the successor of an earlier non-counterexample.
These elementary properties of integers are our starting mathematical
assumptions, not consequences of cryptography.

\subsection{Remainders and positional representations}

\begingroup\mdfsetup{nobreak=true}
\begin{theorem}[Integer division]\label{thm:division}
For an integer \(a\) and positive integer \(b\), unique integers
\(q,r\) satisfy \(a=qb+r\) and \(0\le r<b\).
\end{theorem}
\endgroup
\begin{proof}
Choose the least nonnegative integer in the set of integers \(a-kb\).
Such integers exist by taking \(k\) sufficiently negative. Its value
\(r\) cannot be at least \(b\), because \(r-b\) would be a smaller
nonnegative member. For uniqueness, two decompositions give
\((q-q')b=r'-r\). The right side has absolute value less than \(b\);
the only multiple of \(b\) with that property is zero.
\end{proof}

We write \(a\bmod b\) for this nonnegative remainder, including when
\(a\) is negative. The floor \(\lfloor x\rfloor\) is the greatest
integer not exceeding \(x\); the ceiling \(\lceil x\rceil\) is the
least not smaller. For positive integer \(b\), the quotient above is
\(\lfloor a/b\rfloor\).

\begin{definition}[Digits and their value]\label{def:digits}
For a base, or radix, \(B\ge2\), an \(n\)-digit little-endian sequence
\((d_0,\ldots,d_{n-1})\), with \(0\le d_i<B\), represents
\[
 \val_B(d_0,\ldots,d_{n-1})=\sum_{i=0}^{n-1}d_iB^i.
\]
The sum sign adds one term for each index from zero through \(n-1\).
The highest possible value is \(B^n-1\). Big-endian order reverses the
digits, so its first digit has weight \(B^{n-1}\).
\end{definition}

Here \([0,B^n)\) denotes the integers from zero up to, but excluding,
\(B^n\). Repeated division by \(B\) constructs the digits of every
integer in this interval: take the remainder as the next digit and continue
with the quotient. Uniqueness of division proves uniqueness of every digit,
starting with the least significant. Thus length-\(n\) digit strings
are in bijection with the integers in that interval. A leading zero
does not change the value, but it does change the length; lengths belong
to the encoding specification.

A \emph{bit} has radix two. The notation \(\LE_k(a)\) means the
\(k\)-bit little-endian encoding of an integer \(0\le a<2^k\).
A \emph{byte} is an integer between zero and 255; a byte string is a
sequence of bytes. We write
\(\LE^{\rm byte}_n(a)\) and \(\BE^{\rm byte}_n(a)\) when an
\(n\)-byte order must be explicit. A hexadecimal digit has radix 16;
two such digits specify a byte. For instance \(\mathtt{0x53}\) denotes
the integer \(5\cdot16+3\), not a decimal string containing ``53''.
Concatenation \(s\mathbin\Vert t\) appends the sequence \(t\) to
\(s\); the empty sequence is \(\epsilon\).

\subsection{Canonical encodings and unambiguous tuples}

\begin{definition}[Canonical encoding]\label{def:canonical}
An encoding of a set \(S\) is an injective function
\(\operatorname{enc}:S\to\{0,\ldots,255\}^*\), together with a
decoder that returns \(x\) exactly on \(\operatorname{enc}(x)\) and
returns \(\bot\) on every other string. The star denotes all finite
lengths, not just one fixed length.
\end{definition}

Suppose an integer is required to lie in \([0,p)\) and occupies 32
bytes. The decoder must first recover the integer and check that bound.
Reducing an arbitrary 32-byte integer modulo \(p\) defines a different
operation: it accepts multiple encodings of the same remainder. Rejection
and reduction cannot be interchanged merely because subsequent
arithmetic produces the same value.

Concatenation of variable-length strings is not an injective tuple
encoding: \((\texttt{a},\texttt{bc})\) and
\((\texttt{ab},\texttt{c})\) produce the same concatenation. Fixed
component lengths resolve the ambiguity. Alternatively, encode each
component's length in a specified format before its bytes. The format
must itself be unambiguous and enforce its bounds. A type or domain tag
distinguishes different constructions; it does not recover component
boundaries that were never encoded.

Bitwise exclusive-or, written \(\oplus\), adds corresponding bits
modulo two. It is not ordinary integer addition: carries are discarded.
A word is a fixed-length sequence of bits. For a \(w\)-bit word, rotation
by \(j\) positions moves each bit \(j\) places, wrapping around the ends.
A logical shift inserts zero bits and
discards bits that pass the boundary; a rotation does neither. These
operations, together with addition modulo \(2^w\), specify the integer
operations used later by byte-oriented algorithms.

\section{Prime fields and modular arithmetic}\label{sec:intmod}

Finite arithmetic lets a protocol store every number in a bounded space.
It also changes the meaning of equality. We first construct the
arithmetic, then prove the bounds needed to recover statements about
ordinary integers from statements about residues.

\subsection{Divisibility, greatest common divisors and inversion}

An integer \(d\) divides \(a\), written \(d\mid a\), if
\(a=dk\) for some integer \(k\). A common divisor of \(a,b\) divides
every integer combination \(ua+vb\). For integers not both zero,
\(\gcd(a,b)\) is their greatest positive common divisor.

\begin{theorem}[Euclid and B\'ezout]\label{thm:bezout}
The repeated replacement \((a,b)\leftarrow(b,a\bmod b)\), starting
with nonnegative integers and \(b>0\), terminates. Its last nonzero
remainder is \(g=\gcd(a,b)\). Integers \(u,v\) can be computed with
\(ua+vb=g\).
\end{theorem}
\begin{proof}
Writing \(a=qb+r\) shows that the common divisors of \(a,b\) are
exactly those of \(b,r\). Positive second entries strictly decrease,
so a zero remainder is reached. At that point the positive first entry
is the greatest common divisor. Initially, the two entries have
coefficient pairs \((1,0)\) and \((0,1)\) as combinations of the
original inputs. If the current coefficient pairs are \((u,v)\) and
\((u',v')\), the new remainder has pair
\((u-qu',v-qv')\). Updating these pairs alongside the remainders gives
the claimed identity constructively.
\end{proof}

A positive integer \(p>1\) is prime when its only positive divisors
are 1 and itself. If \(p\nmid a\), then \(\gcd(a,p)=1\), so
Theorem~\ref{thm:bezout} gives \(ua+vp=1\). Consequently, if
\(p\mid ab\) and \(p\nmid a\), multiplying the identity by \(b\)
shows \(p\mid b\). This is the prime-divisor property used below.
Repeatedly factoring a composite integer proves existence of a product
of primes, because each factor is smaller. Applying the prime-divisor
property to two products and cancelling matching factors proves
uniqueness, apart from ordering. This also justifies describing a
positive integer by its prime-power factors.

\subsection{Constructing a field}

We write \(a\equiv b\pmod m\) when \(m\mid(a-b)\). Congruence is
an equivalence relation: \(a\equiv a\); \(a\equiv b\) implies
\(b\equiv a\); and \(a\equiv b\), \(b\equiv c\) imply
\(a\equiv c\), all with the same modulus. These three properties
are called reflexivity, symmetry and transitivity, respectively.
Replacing either operand by a congruent integer leaves its sum and
product congruent to the original result. We may therefore compute
with representatives \(0,\ldots,m-1\), reducing after each operation.
Addition, multiplication and subtraction obey the usual associative,
commutative and distributive laws because their integer counterparts do.

\begin{definition}[Field]\label{def:field}
A field is a set with distinct elements \(0,1\), associative and
commutative addition and multiplication, additive identity zero and
multiplicative identity one, an additive inverse for every element,
and a multiplicative inverse for every nonzero element. Multiplication
distributes over addition.
\end{definition}

\begin{proposition}[Prime fields]\label{prop:prime-field}
For prime \(p\), the residues with arithmetic modulo \(p\) form a
field \(\F_p\). Its inverses are computed by extended Euclid. For
composite \(m>1\), the residues modulo \(m\) do not form a field.
\end{proposition}
\begin{proof}
All laws except inversion have just been inherited from the integers.
For \(0<a<p\), B\'ezout gives \(ua+vp=1\), so \(u\bmod p\)
is an inverse of \(a\). If \(m=ab\) with \(1<a,b<m\), the two
nonzero residues \(a,b\) have product zero. An inverse of \(a\)
would imply \(b=0\), a contradiction.
\end{proof}

In every field, \(ab=0\) implies \(a=0\) or \(b=0\): multiply by
\(a^{-1}\) when \(a\ne0\). We may cancel a common nonzero factor
for the same reason. Division \(a/b\) means \(ab^{-1}\) and is
undefined when \(b=0\). An algorithm that uses a division must prove
the denominator nonzero or specify a failure case.

\begin{example}[A computed inverse]\label{ex:inverse}
The Euclidean chain is
\[
97=5\cdot17+12,\quad17=12+5,\quad12=2\cdot5+2,\quad5=2\cdot2+1.
\]
Back substitution gives \(1=40\cdot17-7\cdot97\). Thus the inverse
of 17 in \(\F_{97}\) is 40. The example checker verifies the final
identity independently of the displayed chain.
\end{example}

\subsection{Efficient powers and simultaneous inverses}

To compute \(a^e\) for a nonnegative integer \(e\), initialise
\(r=1,b=a,k=e\). While \(k>0\), multiply \(r\) by \(b\) if
\(k\) is odd, replace \(b\) by \(b^2\), and replace \(k\) by
\(\lfloor k/2\rfloor\). The invariant \(rb^k=a^e\), interpreted
before each iteration, proves the result. At most twice the bit length
of \(e\) multiplications are needed. A secret-dependent branch in this
description is mathematical pseudocode, not permission for a
timing-dependent implementation.

There is also a useful way to invert many nonzero elements with one
inversion. For \(a_1,\ldots,a_n\ne0\), form prefix products
\(P_0=1\), \(P_i=P_{i-1}a_i\). Invert \(P_n\), obtaining
\(u=P_n^{-1}\). In reverse order, output \(a_i^{-1}=uP_{i-1}\)
and update \(u\leftarrow ua_i\). The update makes
\(u=P_{i-1}^{-1}\), proving each output. Zero inputs must be removed
and handled separately; the product method does not invert zero.

\subsection{When a field equality is an integer equality}

The notation \([a]_p\) denotes the residue of an integer \(a\).
The map \(a\mapsto[a]_p\) preserves sums and products, but it is
not injective on all integers. In particular, there is no order relation
on \(\F_p\) compatible with the usual ordered-field laws. The
representative 96 is larger than 1 as an integer, but in \(\F_{97}\)
it is also \(-1\).

\begin{lemma}[No wraparound]\label{lem:no-wrap}
If integers \(a,b\) satisfy \([a]_p=[b]_p\) and \(|a-b|<p\),
then \(a=b\). More generally, reduction modulo \(p\) is injective
on any interval of integers whose largest and smallest members differ
by less than \(p\).
\end{lemma}
\begin{proof}
The difference is a multiple of \(p\). Every nonzero such multiple
has absolute value at least \(p\).
\end{proof}

\begin{corollary}[A bounded balance equation]\label{cor:balance-bound}
Let \(0\le u_i<U_i\), \(0\le v_j<V_j\), and let \(b\) be a
signed integer with \(|b|\le B\). If
\[
 \sum_i(U_i-1)+\sum_j(V_j-1)+B<p,
\]
then the field equation
\(\sum_i[u_i]_p-\sum_j[v_j]_p=[b]_p\) implies the corresponding
integer equation.
\end{corollary}
\begin{proof}
The absolute value of the difference between the two sides is bounded
by the displayed sum. Apply Lemma~\ref{lem:no-wrap}. Tighter
one-sided bounds are possible; this sufficient bound does not assume
that either side is nonnegative.
\end{proof}

A field constraint \(b(b-1)=0\) forces \(b\in\{0,1\}\).
Therefore constraints of this form together with
\(x=\sum_{i=0}^{k-1}2^ib_i\) describe a \(k\)-bit integer in the
field. When \(2^k\le p\), the represented integer is unique. When
\(2^k>p\), bit constraints alone do not prove that the representation
is canonical modulo \(p\); an additional less-than-\(p\) condition
is necessary. A comparison can be checked by a binary subtraction with
explicit borrow bits, or by a range-constrained difference, provided
its own field-to-integer bound is established.

\begin{example}[A false integer balance]\label{ex:wrap}
The equality \(80+80=63\) holds in \(\F_{97}\), but not in
\(\Z\). All three numbers are themselves canonical field elements.
Canonical inputs do not prevent their sum from wrapping. A proof of a
monetary balance must bound the complete signed difference, not merely
each encoded input.
\end{example}

A field homomorphism is a function between fields preserving one,
addition and multiplication. It necessarily preserves zero, negatives
and inverses: apply it to \(0+0=0\), \(a+(-a)=0\) and
\(aa^{-1}=1\). It is injective, since a nonzero element mapped
to zero would have its inverse product mapped to zero instead of
one. Such an inclusion is also called a field embedding. A bijective
field homomorphism is a field isomorphism; its inverse preserves
the same operations by applying the original map and using
injectivity. These names concern preservation of both operations,
not merely a bijection of the underlying sets.

For distinct primes \(p,r\), conversion from a field element to a
scalar means a specified sequence of operations: take its canonical
integer representative, optionally check an additional bound, then
reduce or canonically encode for the destination. This conversion is
not a field homomorphism. Indeed, a map preserving one and addition
would send \(p\cdot1=0\) in \(\F_p\) to \(p\cdot1\ne0\)
in \(\F_r\). Coordinates and scalars must consequently retain their
types even when their moduli have nearly the same size.

\section{Groups and scalar multiplication}\label{sec:groups}

Field elements support two operations. For the objects called points we
need only one operation and its inverse. The group laws organise that
operation; they do not say that recovering a multiplier from its result
is difficult.

\subsection{The operation and its order}

\begin{definition}[Abelian group]\label{def:group}
An abelian group \((\G,+)\) has an associative and commutative
operation, an identity \(\mathcal O\) satisfying
\(P+\mathcal O=P\), and for each \(P\) an inverse \(-P\) with
\(P+(-P)=\mathcal O\). Its order is its number of elements when
finite. Subtraction means adding an inverse.
\end{definition}

The identity and inverses are unique. If both \(I,J\) are identities,
then \(I=I+J=J\); if both \(Q,R\) invert \(P\), then
\(Q=Q+(P+R)=(Q+P)+R=R\). Cancellation follows by adding the inverse
to both sides of an equality. A subgroup is a subset containing the
identity and closed under addition and inverses; the group laws restrict
to it.

For an integer \(k\ge0\), define \([k]P\) to be the sum of \(k\)
copies of \(P\), with \([0]P=\mathcal O\). For negative \(k\),
set \([k]P=-[-k]P\). Associativity and induction give
\[
 [a+b]P=[a]P+[b]P,\qquad [ab]P=[a]([b]P),
 \qquad [a](P+Q)=[a]P+[a]Q.
\]
The last identity uses commutativity. Double-and-add computes \([k]P\)
by the same binary method as field exponentiation, replacing
multiplication by point addition and squaring by doubling. In a finite
group, some positive multiple of \(P\) is the identity: among more
multiples than group elements, two coincide. The least such positive
integer is the order \(\ord(P)\).

\begin{proposition}[Scalar reduction]\label{prop:scalar-reduction}
If \(P\) has order \(r\), then \([a]P=[b]P\) exactly when
\(a\equiv b\pmod r\). The subgroup generated by \(P\) has exactly
\(r\) elements.
\end{proposition}
\begin{proof}
Divide \(a-b\) by \(r\). Its remainder \(s\) satisfies
\([a-b]P=[s]P\). An identity with \(0<s<r\) would contradict
minimality of \(r\). Conversely, a multiple of \(r\) acts as zero.
The multiples indexed by \(0,\ldots,r-1\) are therefore distinct and
include all multiples.
\end{proof}

\subsection{Counting subgroups and prime-order groups}

For a subgroup \(H\), a coset is the set
\(P+H=\{P+Q:Q\in H\}\). Translation is a bijection from \(H\)
onto this set. Two cosets that share an element are equal: from
\(P+h=Q+h'\) one obtains \(P-Q=h'-h\in H\), which translates
one description into the other. Cosets thus partition a finite group
into equally sized pieces.

\begin{theorem}[Lagrange's counting theorem]\label{thm:lagrange-group}
The order of a subgroup divides the order of a finite group. In
particular, the order of each element divides the group order.
\end{theorem}
\begin{proof}
Count the partition into cosets, then apply the result to the subgroup
generated by the element.
\end{proof}

If the group order is prime \(r\), every nonidentity point has order
\(r\); hence every such point generates the whole group. A group
generated by one element is called cyclic. After choosing a generator
\(G\), the function \(a\mapsto[a]G\) is a bijection from \(\F_r\)
to the group. Multiplication by a nonzero scalar is also a bijection,
with inverse multiplication by its field inverse. These are algebraic
statements even when computing the inverse of the first bijection from
the encoded point is not known to be efficient.

The nonzero elements \(\F_p^\times\) form a multiplicative group of
order \(p-1\). Lagrange therefore proves Fermat's identity
\(a^{p-1}=1\) for \(a\ne0\). The formula
\(a^{-1}=a^{p-2}\) is an alternative inversion algorithm. Unlike
\(\F_p^\times\), a prime-order point group has no nontrivial proper
subgroups. When a larger point group contains a prime-order subgroup
of order \(r\), the ratio of the group order to \(r\) is the
cofactor. Multiplying by a cofactor and checking membership are
different operations; a protocol must specify which it requires.

\subsection{Linear combinations of points}

A map \(f:\G\to\G'\) preserving addition is a homomorphism.
It preserves the identity, inverses and integer scalar multiplication
by the laws already proved. A bijective homomorphism is an isomorphism;
it preserves the abstract operation, not necessarily the representation
or the cost of computing its inverse.

For points \(G_0,\ldots,G_{n-1}\) in a prime-order group and scalars
\(a_i\in\F_r\), the multiscalar multiplication is
\[
 \ip{\mathbf a}{\mathbf G}=\sum_{i=0}^{n-1}[a_i]G_i.
\]
It is linear in each argument where addition and scalar multiplication
are defined. Evaluating each term by double-and-add is a complete
algorithm; faster implementations share work among the terms but do
not change the result. Relative to any fixed generator, every point
has a scalar representation. Thus a list of group generators
always satisfies some relation \(\sum_i[a_i]G_i=\mathcal O\)
with coefficients not all zero when at least two bases are present.
Cryptographic uses of several generators rely instead on a claim that
certain relations between their scalars are computationally unknown.

\section{Probability and computation}\label{sec:probability}

A random choice is specified by a distribution, not by the word
``random''. In particular, a value can be uniform before an observation
and strongly biased after it. We develop the counting and conditioning
rules needed to make such distinctions explicit.

\subsection{Real quantities and approximation}

The rational numbers \(\Q\) are fractions \(a/b\) of integers,
with \(b\ne0\), identifying \(a/b\) and \(c/d\) when
\(ad=bc\). Some limits of increasingly accurate rational
approximations, such as the positive square root of two, are not
rational. We use the real numbers \(\R\) to include such limits.
Concretely, an exact real can be specified by nested closed
intervals with rational endpoints and lengths tending to zero.
The intervals describe approximations at arbitrary accuracy,
not a promise that every real has a finite decimal expansion.

The foundational property used here is completeness: every
nonempty set of reals bounded above has a least upper bound,
its supremum. This is an explicit property of the real number
system, alongside its ordered-field laws; we do not require
the reader to know calculus. The infimum is the corresponding
greatest lower bound. Nested nonempty closed bounded intervals
whose lengths tend to zero have exactly one common point:
take the supremum of their left endpoints, and use the
shrinking lengths for uniqueness. This explains the interval
representation and the convergence of binary search.

The notation \([a,b]\) includes both endpoints; \((a,b)\)
excludes both, with mixed brackets for one included endpoint.
A sequence \(x_n\) tends to \(x\) if, for every accuracy
\(\epsilon>0\), eventually \(|x_n-x|<\epsilon\). A function
is continuous at \(x\) if arbitrarily small changes in its
input near \(x\) make its output arbitrarily close to \(f(x)\):
for every \(\epsilon>0\) there is \(\delta>0\) such that
\(|u-x|<\delta\) implies \(|f(u)-f(x)|<\epsilon\).
Finite sums and products preserve limits, as follows by
bounding the difference of sums or expanding a difference of
products. Quotients preserve them when the limiting denominator
is nonzero. These elementary rules justify the approximations
used below.

\subsection{Experiments, conditioning and expectation}

A finite probability space consists of a set \(\Omega\) of outcomes
and nonnegative weights \(w(\omega)\) summing to one. An event is a
subset \(A\subseteq\Omega\), with probability
\(\Pr[A]=\sum_{\omega\in A}w(\omega)\). In a uniform space all
weights equal \(1/|\Omega|\); this definition requires
\(\Omega\ne\varnothing\). The notation \(X\gets S\) means choose
uniformly from a nonempty finite set \(S\), using fresh independent
randomness unless another distribution is specified.
For a finite distribution, its support is the set
of outcomes with positive probability.

For \(\Pr[B]>0\), conditional probability is
\[
 \Pr[A\mid B]=\frac{\Pr[A\cap B]}{\Pr[B]}.
\]
Events are independent when \(\Pr[A\cap B]=\Pr[A]\Pr[B]\).
For mutually independent trials this factorisation extends to every
selection of events from the different trials; pairwise independence
alone does not establish it. A partition \(B_i\) of the outcome space
gives the total-probability identity
\(\Pr[A]=\sum_i\Pr[B_i]\Pr[A\mid B_i]\), omitting zero-probability
parts. Rearranging the definition gives Bayes' identity
\(\Pr[B\mid A]=\Pr[A\mid B]\Pr[B]/\Pr[A]\).

A random variable is a function on outcomes. For a real-valued finite
random variable, its expectation is the weighted average
\(\E[X]=\sum_\omega w(\omega)X(\omega)\). Rearranging finite sums
proves linearity \(\E[aX+bY]=a\E[X]+b\E[Y]\), without an independence
requirement. An indicator \(1_A\) is one on \(A\) and zero elsewhere;
its expectation is \(\Pr[A]\). The variance
\(\operatorname{Var}(X)=\E[(X-\E X)^2]\) measures squared deviation
from the mean. Expanding the square gives
\(\operatorname{Var}(X)=\E[X^2]-(\E X)^2\).

\begin{lemma}[Basic probability bounds]\label{lem:prob-bounds}
For events \(A_1,\ldots,A_n\),
\(\Pr[\bigcup_i A_i]\le\sum_i\Pr[A_i]\). For nonnegative \(X\)
and \(t>0\), \(\Pr[X\ge t]\le\E[X]/t\). Consequently,
\(\Pr[|X-\E X|\ge t]\le\operatorname{Var}(X)/t^2\).
\end{lemma}
\begin{proof}
Pointwise, \(1_{\cup_iA_i}\le\sum_i1_{A_i}\) and
\(X\ge t1_{X\ge t}\). Taking expectations gives the first two
bounds. Apply the second to \((X-\E X)^2\) for the third.
\end{proof}

These are the union, Markov and Chebyshev bounds, respectively. None
permits replacing a conditional distribution by an unconditional one.
For instance, if a hidden integer is uniform on a set of size \(N\),
then after one incorrect guess is revealed, the next distinct guess
hits with probability \(1/(N-1)\), not \(1/N\).

\subsection{Sampling, collisions and concentrated counts}

For \(q\) independent uniform choices from \(N\) objects, a fixed
pair collides with probability \(1/N\). The union bound over pairs
gives probability at most \(q(q-1)/(2N)\) for any collision. For
\(q\le N\), multiplying the probabilities of avoiding all earlier
choices gives the exact no-collision probability
\[
 \prod_{j=0}^{q-1}\left(1-\frac jN\right).
\]
The pair count is useful even when it exceeds one, but then its upper
bound is uninformative and should be capped at one.

If \(n\) independent trials succeed with probability \(\theta\),
the count \(X\) has binomial distribution
\[
 \Pr[X=k]=\binom{n}{k}\theta^k(1-\theta)^{n-k},\qquad
 \binom{n}{k}=\frac{n!}{k!(n-k)!}.
\]
Here \(n!=1\cdot2\cdots n\), with \(0!=1\). The coefficient
counts the choices of the \(k\) successful positions. Summing over all
subsets proves that the probabilities sum to one. Writing the count as
a sum of indicators gives \(\E X=n\theta\); independence makes the
cross terms in its variance vanish, giving
\(\operatorname{Var}(X)=n\theta(1-\theta)\).

Without replacement, sampling \(n\) distinct objects uniformly from
\(N\) objects, of which \(K\) are marked, instead gives
\[
 \Pr[X=k]=\frac{\binom{K}{k}\binom{N-K}{n-k}}{\binom{N}{n}}.
\]
Impossible choices have coefficient zero. Count the marked and unmarked
subsets to obtain this formula. These two distributions cannot be
exchanged simply because the sample size is small relative to the
population; an approximation needs its own error bound.

For sharper upper tails, apply Markov to \(t^X\), with \(t>1\).
Let \(X\) be a sum of independent zero/one variables, where
variable \(i\) equals one with probability \(\theta_i\).
For an integer threshold \(a\), they satisfy the finite bound
\begin{equation}\label{eq:tilted-tail}
 \Pr[X\ge a]\le t^{-a}\prod_i(1-\theta_i+\theta_i t).
\end{equation}
Indeed, the expectation of a product of independent variables factors,
and \(\E[t^{1_A}]=1-\theta+\theta t\). In the equal-probability
case, for \(0<\theta<1\) and \(n\theta<a<n\), substituting
\(t=a(1-\theta)/(\theta(n-a))>1\) yields
\[
 \Pr[X\ge a]\le
 \left(\frac{n\theta}{a}\right)^a
 \left(\frac{n(1-\theta)}{n-a}\right)^{n-a}.
\]
This explicit bound needs no replacement distribution and no unstated
limit. If \(\theta=0\), then \(X=0\) certainly and every positive
threshold has probability zero; if \(\theta=1\), then \(X=n\)
certainly. Related sampling models must justify their own expectation
factorisation or another bound replacing it.

\subsection{Countable experiments and waiting times}

An algorithm may retry indefinitely. A finite probability space no
longer describes all its possible running times. For a countable set of
outcomes, meaning outcomes that can be listed by natural-number
indices, we assign nonnegative weights with sum one; an infinite sum of
nonnegative terms is the supremum of its finite partial sums. This
definition permits rearranging such terms, since both arrangements
have the same finite subsums. Probabilities of increasing events tend
to the probability of their union by the same reasoning. We need no
general integration theory for these uses.

For an integer-valued \(T\ge0\), write
\(T=\sum_{j\ge0}1_{T>j}\) pointwise. Nonnegative rearrangement gives
the tail-sum formula
\begin{equation}\label{eq:tail-sum}
 \E T=\sum_{j\ge0}\Pr[T>j],
\end{equation}
allowing an infinite expectation. With independent success probability
\(\theta>0\) on every trial, the number \(T\ge1\) of trials up to
and including the first success has
\(\Pr[T>j]=(1-\theta)^j\) and \(\E T=1/\theta\). The geometric
series identity follows by summing its first \(m\) terms and taking
the limit. Success occurs eventually with probability one, but that
statement by itself would not imply finite expected time.

Waiting for \(r\ge1\) successes requires a number of failures \(K\)
with
\[
 \Pr[K=k]=\binom{k+r-1}{k}\theta^r(1-\theta)^k.
\]
The final trial must succeed; choose the \(k\) failures among the
preceding \(k+r-1\) positions. Equivalently, the total waiting time is
a sum of \(r\) geometric waiting times, so it has expectation
\(r/\theta\). This is the negative-binomial waiting model. It differs
from fixing an elapsed time first and counting successes during it.

\subsection{Integrals and continuous distributions}

For a bounded function \(f\) on \([a,b]\), divide the interval
at finitely many points \(a=x_0<\cdots<x_n=b\). Multiplying
each subinterval's width by the infimum of \(f\) there gives
a lower sum; using its supremum gives an upper sum. Refining
a partition cannot decrease the lower sum or increase the
upper sum. If the supremum of all lower sums equals the
infimum of all upper sums, their common value is the integral
\(\int_a^b f(x)\,dx\). It is thus a quantity constructed
by progressively tighter finite arithmetic bounds.
Choosing one point in each subinterval instead gives a Riemann
sum, the sum of its function value times that subinterval's
width. It lies between the lower and upper sums. The mesh is
the largest width in the partition.

Continuous functions on a closed bounded interval have such
an integral. Here are the needed details. Repeated interval
bisection shows that any sequence in \([a,b]\) has a
subsequence converging to a point in that interval: retain
a half containing infinitely many sequence entries at each
stage and choose entries with increasing indices. If a
continuous function were unbounded, a sequence on which its
absolute value grows beyond every bound would have such a
subsequence, contradicting continuity at its limiting point.
Likewise, if there were no uniform input tolerance for a
fixed output accuracy, there would be pairs \(u_n,v_n\)
with \(|u_n-v_n|\to0\) but a fixed positive separation
of their outputs. A converging subsequence of \(u_n\),
and hence of \(v_n\), contradicts continuity. Thus for
any \(\epsilon>0\) we can choose a mesh small enough that
the upper and lower values in each subinterval differ by
less than \(\epsilon\). The upper and lower sums differ
by at most \(\epsilon(b-a)\), proving integrability.

The finite sums prove positivity, additivity over adjoining
intervals, and linearity in the integrand. They also prove
\(|\int_a^b f|\le(b-a)\sup|f|\). Define an integral with
reversed endpoints by \(\int_b^a f=-\int_a^b f\).
A shift and rescaling of the variable, \(t=c+du\), rescales every
partition width by \(|d|\); accounting for endpoint order
gives
\[
 \int_\alpha^\beta f(c+du)d\,du
 =\int_{c+d\alpha}^{c+d\beta}f(t)\,dt.
\]
This substitution identity is proved directly from the
partition sums, without assuming a differentiation rule.

A continuous probability density on an interval is a
nonnegative function \(f\) whose integral over its domain
is one. An interval probability is its integral on that
interval. For an unbounded domain, define the total integral
as the increasing limit over bounded intervals. A cumulative
distribution function, abbreviated CDF, records
\(F(x)=\Pr[X\le x]\). For a density on \([a,b]\),
it is \(F(x)=\int_a^x f(t)\,dt\) inside the interval,
zero below it and one above it. Individual points have
probability zero; probability is spread over intervals,
not assigned a positive weight to each real point.
For a density \(f\), the expectation of a function \(g(X)\)
is \(\int g(x)f(x)\,dx\) when the required integrals exist
and the positive and negative parts are not both infinite.
Our applications use continuous or piecewise-continuous
integrands and their stated nonnegative limits. Piecewise-continuous
means continuous on finitely many adjoining pieces of each bounded
integration interval, with finite one-sided endpoint limits; the integral
construction does not make every arbitrary bounded function
integrable. For nonnegative
quantities, rearrangement of nonnegative partition sums gives
the continuous counterpart of \eqref{eq:tail-sum}:
\[
 \E X=\int_0^\infty\Pr[X>t]\,dt\qquad(X\ge0).
\]
Indeed, approximate each value \(x\) from below by a sum
of equally spaced threshold indicators, apply linearity
to the finite sum, and take increasingly fine meshes and
increasing cutoffs. Both sides are limits of the same
nonnegative finite sums.

An ideal uniform real \(U\) on \([0,1]\) can be represented
by an infinite sequence of independent fair bits:
\(U=\sum_{i\ge1}b_i2^{-i}\). Fixing the first \(k\)
bits selects an interval of length \(2^{-k}\) with that
probability. Approximating arbitrary interval endpoints by
such binary intervals gives probability equal to interval
length. Endpoints with two binary expansions have probability
zero and do not change the distribution. This is an ideal
continuous experiment, not an algorithm that can return
infinitely many bits in finite time.

\begin{lemma}[Sampling by the inverse CDF]\label{lem:inverse-cdf}
If \(F\) is continuous and strictly increasing from zero
to one on an interval, then \(X=F^{-1}(U)\)
has CDF \(F\), where \(U\) is uniform on \([0,1]\).
Endpoints or infinite endpoints can be defined by limits.
\end{lemma}
\begin{proof}
Monotonicity gives
\(F^{-1}(U)\le x\) exactly when \(U\le F(x)\).
The latter event has probability \(F(x)\). Existence of
the inverse follows by binary search and continuity,
and uniqueness follows by strict monotonicity.
\end{proof}

A finite implementation may replace \(U\) by a uniform
grid point \(j/2^k\), \(0\le j<2^k\). For every
threshold, the resulting interval count differs from its
uniform-real probability by at most \(2^{-k}\). This
is a bound on CDF error, not closeness for every possible
event. The event that the output lies among the finitely
many values the discrete sampler can return
has probability one in the discrete experiment and zero
in the continuous experiment. A protocol requiring a discrete
distribution must specify its actual rounding or rejection
algorithm rather than claiming that floating-point inversion
is exact continuous sampling.

\subsection{Natural logarithms, exponentials and their identities}

For \(x>0\), define the natural logarithm constructively by
\[
 \ln x=\int_1^x\frac{dt}{t}.
\]
The integrand is continuous on every positive closed bounded
interval, so the preceding construction applies. Positivity
of \(1/t\) makes \(\ln\) strictly increasing. The integral
bound on a short interval proves its continuity. Splitting
the interval at \(a>0\) and substituting \(t=au\) proves
\[
 \ln(ab)=\ln a+\ln b,\qquad
 \ln(1/a)=-\ln a,\qquad
 \ln(a/b)=\ln a-\ln b.
\]
For example, \(\ln2\ge1/2>0\), so
\(\ln(2^n)=n\ln2\) grows without bound and
\(\ln(2^{-n})=-n\ln2\) decreases without bound.
Continuity and binary search therefore give, for every
real \(y\), a unique positive real \(x\) with \(\ln x=y\).
Define \(\exp(y)\) to be that real, and write \(e=\exp(1)\).
The product law above gives
\[
 \exp(u+v)=\exp(u)\exp(v),\qquad
 \exp(-u)=1/\exp(u).
\]
For \(a>0\) and real \(u\), define
\(a^u=\exp(u\ln a)\). This agrees with integer powers
and supplies the exact meaning of fractional powers in
sampling formulas. Logarithm to base \(a\ne1\) is
\(\log_a x=\ln x/\ln a\), the inverse of \(u\mapsto a^u\).

The exponential also has the convergent series
\begin{equation}\label{eq:exp-series}
 \exp(x)=\sum_{k=0}^{\infty}\frac{x^k}{k!}.
\end{equation}
To justify the identity without presupposing calculus,
call the series \(E(x)\). Absolute convergence follows
because successive absolute terms eventually have ratio
at most one half. Multiplying two absolutely convergent
series and grouping by total degree, then using the finite
binomial identity, gives \(E(x+y)=E(x)E(y)\). These
regroupings are legitimate: absolute tail sums bound the
discarded terms. In particular \(E(x)E(-x)=1\).
The series is positive for \(x\ge0\), so it is positive
everywhere. For \(|h|\le1\), its terms of degree at
least two have total absolute value at most a fixed
constant times \(h^2\). Hence
\((E(h)-1)/h\to1\), and the product identity makes
\(E\) continuous everywhere.

The continuous function \(L(x)=\ln E(x)\) is additive.
Every continuous additive function equals \(cx\) for
a fixed constant \(c\): prove it first at integers,
then rational fractions, then pass to real limits.
For \(v\to0\), the integral bounds for \(1/t\)
between 1 and \(1+v\) give
\(\ln(1+v)/v\to1\). Substituting \(v=E(h)-1\)
therefore gives \(L(h)/h\to1\), forcing \(c=1\).
Thus \(\ln E(x)=x\), which proves
\eqref{eq:exp-series} from the inverse definition.

The shift and rescaling \(t=1-x\) now yield the exact
integral identity needed for inverse-distance sampling:
\begin{equation}\label{eq:log-integral}
 \int_a^b\frac{dx}{1-x}
 =\ln(1-a)-\ln(1-b),\qquad a<b<1.
\end{equation}
The equivalent notation \(\int dx/(1-x)=-\ln(1-x)+C\)
means that differences of the expression at two endpoints
give the definite integral; the arbitrary constant cancels.
No knowledge of integration notation is assumed before
this construction.

An exponential waiting time of rate \(\lambda>0\) has
CDF \(F(t)=1-\exp(-\lambda t)\) for \(t\ge0\),
and zero for \(t<0\). It is sampled by
\(-\ln(1-U)/\lambda\); the event \(U=1\) has probability
zero and is handled by a limiting endpoint. Its survival
probability is \(\exp(-\lambda t)\), so
\[
 \Pr[T>s+t\mid T>s]=\exp(-\lambda t)=\Pr[T>t].
\]
This is the memoryless property, derived from the exponential
product identity, not assumed from the phrase ``waiting time''.
Its density is \(\lambda\exp(-\lambda t)\). One can
verify the corresponding integral directly: the series gives
\(\exp(-\lambda h)=1-\lambda h+R(h)\), where
\(|R(h)|\le Ch^2\) for a fixed constant on bounded
intervals. Summing the increments of
\(1-\exp(-\lambda t)\) along a partition differs from
the density's Riemann sum by at most a constant times
the largest partition width, which tends to zero.
The tail-integral formula and this density calculation give
\(\E T=1/\lambda\).

For completeness, a Poisson distribution with parameter \(\lambda>0\)
assigns probability \(e^{-\lambda}\lambda^k/k!\) to \(k\ge0\).
One may define its normalising constant directly by
\(e^\lambda=\sum_{k\ge0}\lambda^k/k!\); the series converges because
the ratio of consecutive terms eventually becomes less than one half.
Reindexing nonnegative sums gives mean and variance \(\lambda\).
A counting process records the number \(N(t)\) of events
up to each time \(t\ge0\), starting from zero. An increment
\(N(t+h)-N(t)\) counts events in the subsequent interval.
Describing this as a Poisson process additionally assumes that
increments on disjoint intervals are independent and have this
distribution with \(\lambda\) proportional to
the interval length. That is a modelling hypothesis about event
arrivals, not a consequence of the binomial formula or of a hash
function's definition.

\subsection{Statistical distance and coupled executions}

\begingroup\mdfsetup{nobreak=true}
\begin{definition}[Statistical distance]\label{def:stat-distance}
For distributions \(P,Q\) on the same finite or countable set,
\[
 \SD(P,Q)=\frac12\sum_x|P(x)-Q(x)|.
\]
\end{definition}
\endgroup

Let \(D=\{x:P(x)>Q(x)\}\). Since both distributions sum to one,
the total positive excess equals the total negative deficit, so
\(\SD(P,Q)=P(D)-Q(D)\). For any event \(A\), the difference
\(P(A)-Q(A)\) cannot exceed that positive excess; reversing the roles
gives \(|P(A)-Q(A)|\le\SD(P,Q)\). Thus distance is exactly the
largest difference in probability of an observable event.

The triangle inequality for absolute values proves
\(\SD(P,R)\le\SD(P,Q)+\SD(Q,R)\). Applying the same deterministic
or randomised computation to both inputs cannot increase their distance:
for each output event, its input-dependent probability is between zero
and one, and summing against \(P-Q\) is bounded by the total positive
excess. This is the data-processing property.

\begin{lemma}[Identical until a bad event]\label{lem:coupling}
Suppose two executions are realised on one probability space and their
observable outputs agree except on an event \(B\). The statistical
distance between their output distributions is at most \(\Pr[B]\).
\end{lemma}
\begin{proof}
For any output event \(A\), its two indicators differ only when the
outputs differ, and their absolute difference is at most one. Their
expectations therefore differ by at most \(\Pr[B]\). Use the event
characterisation of distance.
\end{proof}

A shared realisation of two distributions is a coupling. Its usefulness
is that it can preserve a complete adaptive transcript until a precise
failure occurs. The lemma does not assert that the two executions have
their original distributions \emph{conditioned on} avoiding failure.
Conditioning changes distributions and must be analysed separately.

\subsection{Uniform maps, rejection and reduction bias}

A bijection carries the uniform distribution to uniform. More
generally, a function with exactly \(c\) preimages for every output
carries uniform input to uniform output: the probability of each output
is \(c\) divided by the input-set size. These facts need no
computational assumption.

To sample uniformly from \(\{0,\ldots,m-1\}\), choose
\(k\) with \(2^k\ge m\), sample \(k\) independent bits, and
reject values at least \(m\). Each accepted value had the same
probability before conditioning, so the result is exactly uniform.
The expected number of trials is \(2^k/m\). If instead one reduces
the \(k\)-bit integer modulo \(m\), write \(2^k=qm+s\),
\(0\le s<m\). The first \(s\) residues have \(q+1\) preimages,
and the others have \(q\). Direct substitution in the distance formula
gives
\begin{equation}\label{eq:reduction-bias}
 \SD(U_k\bmod m,U_m)=\frac{s(m-s)}{m2^k}
 \le\frac{m}{4\cdot2^k}.
\end{equation}
The inequality uses \(s(m-s)\le m^2/4\), obtained by completing the
square. A long reduction input makes the bias small; it does not make
it identically zero.

\begin{example}[Visible modular bias]\label{ex:mod-bias}
Reducing a uniform five-bit integer modulo seven produces preimage
counts \((5,5,5,5,4,4,4)\). Its distance from uniform modulo seven
is \(3/56\). The checker counts the preimages and evaluates the
distance using integers, so this example does not rely on floating-point
rounding.
\end{example}

\subsection{Algorithms, costs and security parameters}

An algorithm is a finite description of operations on finite inputs,
with a specified result or failure behaviour. A deterministic algorithm
has no random choices; a randomised one also reads random bits. Its
input length counts bits, not the numerical value of an encoded integer.
Testing every candidate below a \(k\)-bit modulus takes on the order
of \(2^k\) iterations, even though the modulus occupies only \(k\)
bits. Counting a field multiplication as one operation is an algebraic
cost model; a bit-cost claim must additionally account for the size and
implementation of the multiplication.

For nonnegative functions on positive integers, \(f=O(g)\) means
there are constants \(C,n_0\) with \(f(n)\le Cg(n)\) for
\(n\ge n_0\). The notation \(\Omega(g)\) reverses the bound, and
\(\Theta(g)\) asserts both. Logarithms \(\log_2 n\) measure binary
length up to rounding. Polynomial time means a bound \(Cn^d\) for
fixed constants \(C,d\), not constants chosen separately for each
input. A probabilistic polynomial-time algorithm, abbreviated PPT,
uses random bits and obeys such a time bound on every execution.
Expected polynomial time is a different condition and must be named
when used.

A security parameter \(\lambda\) indexes a family of instances and
algorithms. A nonnegative function \(\epsilon(\lambda)\) is negligible
if, for every positive integer \(d\), eventually
\(\epsilon(\lambda)<\lambda^{-d}\). A fixed inverse polynomial is
not negligible. Exponential decay such as \(2^{-\lambda}\) is
negligible: for a fixed \(d\), the ratio of consecutive terms of
\(\lambda^d/2^\lambda\) tends to one half, so eventually the sequence
decreases geometrically to zero. Finite sums of negligible functions,
and their products by fixed polynomials, are negligible by choosing
a stronger exponent in the definition.

An asymptotic claim does not supply a concrete attack cost for one
fixed curve. A concrete statement records running time, query counts
and success probability at the chosen parameters. Likewise, expected
time bounded by \(T\) can be truncated after \(T/\eta\) steps with
abort probability at most \(\eta\), by Markov; truncation is a change
to the algorithm whose failure probability belongs in the final bound.

\section{Linear algebra and uniform masks}\label{sec:linalg}

Linear algebra studies operations that preserve sums and multiplication
by field elements. It supplies two different tools: efficient identities
for rearranging long computations, and exact counting statements about
which observations remain hidden by random masks.

\subsection{Vectors, subspaces and bases}

For a field \(\F\), a vector in \(\F^n\) is an ordered tuple
\(\mathbf v=(v_0,\ldots,v_{n-1})\). Addition and scalar
multiplication operate componentwise. A subspace \(V\subseteq\F^n\)
contains zero and is closed under those operations. A linear combination
of vectors \(\mathbf v_i\) is a sum \(\sum_i a_i\mathbf v_i\).
Their span is the subspace of all their linear combinations. The vectors
are linearly independent if \(\sum_i a_i\mathbf v_i=0\) forces every
coefficient to be zero. A basis is an independent spanning list.

The standard coordinate vectors \(\mathbf e_i\), with a single one
in position \(i\), form a basis of \(\F^n\). Relative to any basis,
the coordinates of a vector are unique: subtracting two expressions
would give a nonzero dependence. If a list is dependent, one vector
whose coefficient in a dependence is nonzero can be expressed using
the others and removed without changing the span. Removing such vectors
from a finite spanning list constructs a basis.

\begin{lemma}[Basis exchange]\label{lem:basis-exchange}
If an independent list has \(k\) vectors and a spanning list has
\(m\) vectors in the same space, then \(k\le m\). All bases of a
space with a finite basis consequently have the same length.
\end{lemma}
\begin{proof}
Express the first independent vector using the spanning list. Some
coefficient is nonzero, so solve for that spanning vector and replace
it by the independent vector, preserving the span. Repeat with the
next independent vector. Its expression must use at least one
not-yet-replaced vector with a nonzero coefficient; otherwise it would
lie in the span of earlier independent vectors. Each step consumes a
different member of the original spanning list, so there can be no
more than \(m\) steps. Apply this in both directions to two bases.
\end{proof}

The common basis length is the dimension. An independent list can be
extended to a basis by adding vectors from a known spanning list when
they increase the span. Thus every subspace of \(\F^n\) has a basis
of at most \(n\) vectors: start with an empty list and add a vector
outside its span until none remains. More than \(n\) additions would
contradict Lemma~\ref{lem:basis-exchange} applied in \(\F^n\).

\subsection{Linear maps, matrices and elimination}

A linear map \(L:V\to W\) satisfies
\(L(a\mathbf v+b\mathbf w)=aL(\mathbf v)+bL(\mathbf w)\).
It is determined by its values on a basis. In coordinate spaces these
values become the columns of a matrix \(A\); the output coordinates
are
\[
 (A\mathbf x)_i=\sum_j A_{ij}x_j.
\]
Matrix multiplication represents composition:
\((AB)_{ij}=\sum_k A_{ik}B_{kj}\). Interchanging finite sums verifies
\((AB)\mathbf x=A(B\mathbf x)\). The identity matrix has ones on
its diagonal and zeros elsewhere.

The kernel \(\ker L\) is the subspace mapped to zero. The image is
a subspace of the codomain; its dimension is the rank. A square matrix
is invertible when its linear map is bijective. Its inverse represents
the inverse map, so multiplication in either order gives the identity.

One can solve \(A\mathbf x=\mathbf b\) constructively by row
operations on the augmented matrix \((A\mid\mathbf b)\). Swap two
rows, multiply a row by a nonzero field element, or add a multiple of
one row to another. Each operation has an inverse of the same form
and preserves the set of solutions. Working from left to right, choose
a nonzero entry as a pivot, move its row into position, scale it to one,
and eliminate entries below it. A row with zero coefficients and a
nonzero right side proves inconsistency. Otherwise, choose arbitrary
values for the nonpivot variables and recover the pivot variables by
back substitution. Eliminating above pivots as well gives an explicit
description without back substitution.

\begin{theorem}[Rank and nullity]\label{thm:rank-nullity}
For a linear map \(L:V\to W\) on a finite-dimensional domain,
\[
 \dim V=\dim\ker L+\dim\im L.
\]
\end{theorem}
\begin{proof}
Take a basis \(\mathbf k_1,\ldots,\mathbf k_s\) of the kernel and
extend it to a basis of \(V\) by
\(\mathbf v_1,\ldots,\mathbf v_t\). The images of these additional
vectors span the image because the kernel basis maps to zero. They
are independent: a combination with zero image would lie in the
kernel and therefore express a dependence in the extended basis.
Their number is the rank, giving the equality.
\end{proof}

In particular, a square matrix is invertible exactly when its kernel
is zero. Elimination tests this and constructs the inverse by applying
the row operations to \((A\mid I)\). For an \(n\)-dimensional
domain over a field with \(q\) elements, a rank-\(r\) linear map
has a kernel with \(q^{n-r}\) elements, since each basis coordinate
has \(q\) choices. Every nonempty fibre
\(L^{-1}(\mathbf y)\) is a translate of that kernel: if
\(L(\mathbf x_0)=\mathbf y\), then
\(L(\mathbf x)=\mathbf y\) precisely when
\(\mathbf x-\mathbf x_0\in\ker L\).

\subsection{Inner products and point-valued linear maps}

The field inner product is
\(\ip{\mathbf a}{\mathbf b}=\sum_i a_i b_i\). It is bilinear,
meaning linear in either argument with the other fixed. It is also
nondegenerate: if its value is zero for every \(\mathbf b\), test
\(\mathbf b=\mathbf e_i\) to obtain \(a_i=0\) for every \(i\).
Unlike a sum of squares over the real numbers, this field-valued
expression need not be nonzero on a nonzero vector paired with itself.
For instance, a nonzero vector can have
\(\ip{\mathbf a}{\mathbf a}=0\) over a finite field. No argument
may infer that a sum of field squares vanishes only when every term
vanishes.

The mixed inner product \(\ip{\mathbf a}{\mathbf G}\) of
\S\ref{sec:groups} has a point as output. The same distribution laws
permit splitting it into halves or changing coordinates. For an
invertible field matrix \(M\), transposition exchanges row and column
indices, and
\[
 \ip{M\mathbf a}{M^{-T}\mathbf b}=\ip{\mathbf a}{\mathbf b},
 \qquad M^{-T}=(M^{-1})^T.
\]
Expanding the products and using \(M^TM^{-T}=I\) proves this
identity. Point-valued versions are proved by the same finite sum
expansion, using the scalar laws in the group. An implementation may
use such identities to combine many checks, but the identities alone
say nothing about which coefficients an adversary knows.

\subsection{What a random linear mask hides}

An affine map is a linear map followed by a translation:
\(\mathbf x\mapsto\mathbf c+L(\mathbf x)\). A translate
\(\mathbf c+V=\{\mathbf c+\mathbf v:\mathbf v\in V\}\)
of a subspace is an affine set. It need not contain zero
and need not be a subspace. These definitions describe
the exact output sets whose probabilities we now count.

\begin{theorem}[Uniform image of an affine map]\label{thm:affine-mask}
Let \(R\) be uniform on \(\F_q^n\), let
\(L:\F_q^n\to\F_q^k\) have rank \(r\), and fix
\(\mathbf c\in\F_q^k\). The vector \(\mathbf c+L(R)\) is
uniform on the affine set \(\mathbf c+\im L\), which has
\(q^r\) elements. It is uniform on the entire codomain exactly when
\(r=k\).
\end{theorem}
\begin{proof}
Every output in the image has \(q^{n-r}\) preimages by the fibre
count above. Its probability is \(q^{n-r}/q^n=q^{-r}\).
Translation by \(\mathbf c\) is a bijection.
\end{proof}

If a secret random variable \(W\) is independent of \(R\) and
\(L\) is surjective, then \(C(W)+L(R)\) has the same uniform
distribution for every possible value of \(W\). Hence the output is
independent of the secret. If the rank is smaller, the affine support
may depend on the secret and reveal information even though every
individual coordinate is uniform.

\begin{example}[Marginal hiding is not joint hiding]
\label{ex:correlated-mask}
Over \(\F_5\), let \(u,v\) be independent uniform values. The map
\((u,v)\mapsto(u+v,u+2v)\) is bijective: subtracting the coordinates
recovers \(v\), then \(u\). Adding any secret pair therefore leaves
the output uniform on all 25 pairs. In contrast,
\((u,v)\mapsto(u+v,2u+2v)\) has only five outputs. After adding a
secret pair \((c_1,c_2)\), the observations satisfy
\(y_2-2y_1=c_2-2c_1\). Each observation separately is uniform, yet
together they reveal this secret linear combination. The checker
enumerates both distributions exactly.
\end{example}

Previously disclosed linear information consumes masking dimensions.
This can be quantified without an independence guess.

\begin{lemma}[Masks after linear observations]\label{lem:conditional-mask}
Let \(R\) be uniform on \(\F_q^n\). Condition on
\(BR=\mathbf b\), for a value in the image of \(B\). Then
\(AR\) is uniform on a translate of \(A(\ker B)\), and
\[
 \dim A(\ker B)=
 \rank\begin{pmatrix}B\\A\end{pmatrix}-\rank B.
\]
In particular, \(k\) output coordinates of \(AR\) remain jointly
uniform exactly when the displayed rank difference equals \(k\).
\end{lemma}
\begin{proof}
Conditioning a uniform input on a fibre makes it uniform on that
fibre, namely \(\mathbf r_0+\ker B\). Apply
Theorem~\ref{thm:affine-mask} to the restriction of \(A\) to
\(\ker B\). Rank--nullity identifies its rank as
\(\dim\ker B-\dim(\ker B\cap\ker A)\). Applying rank--nullity
again to \(B\) and the stacked map proves the formula.
\end{proof}

This lemma has fixed linear maps. If challenges or maps depend
adaptively on previous observations, one must justify the conditional
distribution at each step, including which fresh mask coordinates
remain. Merely counting random coefficients in the original witness
does not prove that a complete adaptive transcript is hidden.

\section{Polynomials and evaluation}\label{sec:polynomials}

A polynomial turns a list of coefficients into a rule for obtaining
field values. Its crucial property is that a nonzero polynomial of
small degree cannot vanish at many different points. That fact supports
both reconstruction from evaluations and probabilistic equality checks.

\subsection{Formal polynomials and division}

A polynomial over \(\F\) is a finite formal sum
\(f(X)=\sum_{i=0}^d a_iX^i\). The symbol \(X\) is an
indeterminate, not a sampled field element. Equality of polynomials
means equality of all coefficients after appending zero coefficients
as necessary. Addition is coefficientwise and multiplication is
convolution:
\[
 (fg)_k=\sum_{i+j=k}a_i b_j.
\]
Associativity, commutativity and distributivity follow by rearranging
finite sums. The set of polynomials is written \(\F[X]\).
Evaluation substitutes \(x\in\F\) for \(X\), giving
\(f(x)=\sum_i a_ix^i\). Evaluation preserves addition and
multiplication, again by expansion.

For nonzero \(f\), its degree is the highest index with nonzero
coefficient. The zero polynomial has degree \(-\infty\), a convention
that makes inequalities about upper degree bounds include zero.
We never use an expression such as a negative number of roots for this
special degree. For nonzero polynomials, multiplying the leading
coefficients gives a nonzero leading coefficient, hence
\(\deg(fg)=\deg f+\deg g\). Also
\(\deg(f+g)\le\max(\deg f,\deg g)\), with possible strict
inequality when leading terms cancel.

\begin{theorem}[Polynomial division]\label{thm:poly-division}
For \(f,g\in\F[X]\) with \(g\ne0\), unique polynomials \(q,r\)
satisfy \(f=qg+r\) and either \(r=0\) or \(\deg r<\deg g\).
\end{theorem}
\begin{proof}
If the current remainder has degree at least \(\deg g\), subtract
the unique scalar multiple of a power of \(X\) times \(g\) that
cancels its leading term. This uses the inverse of the leading
coefficient of \(g\). Its degree strictly decreases, so the process
terminates, and the subtracted multiples give \(q\). For uniqueness,
two outputs would give \((q-q')g=r'-r\). A nonzero left side has
degree at least \(\deg g\); a nonzero right side has smaller degree.
They can therefore only both be zero.
\end{proof}

Evaluation can be performed by Horner's rule. Starting from the highest
coefficient, repeatedly multiply the accumulated value by \(x\) and
add the next coefficient. Expanding the nested expression proves it
equals \(f(x)\), using \(d\) multiplications and \(d\) additions
for a degree-\(d\) polynomial. Polynomial division, likewise, is an
explicit finite algorithm, not a claim that a quotient happens to exist.

\subsection{Roots, multiplicity and equality tests}

\begin{theorem}[Factor and root theorems]\label{thm:root-bound}
For \(a\in\F\), the equality \(f(a)=0\) holds exactly when
\(X-a\) divides \(f\). A nonzero degree-\(d\) polynomial has at
most \(d\) distinct roots in \(\F\).
\end{theorem}
\begin{proof}
Division by \(X-a\) gives a constant remainder; evaluating at \(a\)
identifies it with \(f(a)\). For the root bound, a nonzero constant
has no roots. If a positive-degree polynomial has a root \(a\), write
\(f=(X-a)g\). Every other root is a root of \(g\), since its
\(X-a\) factor is nonzero there. Induction on the degree proves the
bound.
\end{proof}

A root has multiplicity \(m\) if \((X-a)^m\) divides \(f\) but
\((X-a)^{m+1}\) does not. We use this definition only for nonzero
polynomials. Removing factors successively shows that the sum of
multiplicities of distinct roots is at most the degree. The formal
derivative is \(f'(X)=\sum_{i\ge1}i a_iX^{i-1}\), where the
integer coefficient is interpreted in the field. Direct multiplication
proves the product rule \((fg)'=f'g+fg'\). Thus a root of
multiplicity at least two is also a root of \(f'\). Conversely, if
\(f=(X-a)g\), then \(f'(a)=g(a)\); if this vanishes, \(X-a\)
divides \(g\), giving multiplicity at least two. These statements
use a formal operation and require no calculus.

Over \(\F_p\), the nonzero polynomial \(X^p-X\) evaluates to zero
at every field element, by Fermat and the separate zero case. Therefore
polynomials and the functions they induce on a finite field are not
the same objects. Degree restrictions are essential to uniqueness
arguments.

\begin{corollary}[A random evaluation test]\label{cor:random-evaluation}
Fix different polynomials \(f,g\) of degree at most \(d\), and then
sample \(x\) uniformly from a nonempty finite set \(S\subseteq\F\).
The probability that \(f(x)=g(x)\) is at most
\(\min(1,d/|S|)\).
\end{corollary}
\begin{proof}
Apply the root bound to the nonzero difference.
\end{proof}

The word ``then'' is a hypothesis. If the purported polynomial can be
chosen after observing \(x\), the adversary can choose \(X-x\),
which passes at that point with certainty. Any use of the corollary in
a protocol must explain how its polynomial is fixed before the sampled
challenge, or invoke a separate argument that supplies that property.

\subsection{Interpolation and vanishing polynomials}

Let \(H=\{x_0,\ldots,x_{n-1}\}\) contain \(n\) distinct field
elements, with \(n\ge1\). Define
\[
 Z_H(X)=\prod_{i=0}^{n-1}(X-x_i),\qquad
 L_i(X)=\prod_{j\ne i}\frac{X-x_j}{x_i-x_j}.
\]
All denominators are nonzero. The polynomials \(L_i\) have degree
less than \(n\) and satisfy \(L_i(x_j)=1\) when \(i=j\), zero
otherwise.

\begin{theorem}[Interpolation]\label{thm:interpolation}
For any \(y_0,\ldots,y_{n-1}\in\F\), exactly one polynomial of
degree less than \(n\) takes those values on the corresponding points.
It is \(f(X)=\sum_i y_iL_i(X)\).
\end{theorem}
\begin{proof}
The evaluation property of the \(L_i\) proves existence. The
difference of two candidates would have \(n\) distinct roots and
degree less than \(n\), so the root bound forces it to be zero.
\end{proof}

The space of degree-less-than-\(n\) polynomials has coefficient
basis \(1,X,\ldots,X^{n-1}\). Interpolation proves it also has
basis \(L_0,\ldots,L_{n-1}\): these polynomials span by the formula
and are independent by evaluation at each \(x_i\). Evaluation is
therefore an invertible linear change of coordinates. Its matrix is
\(V_{ij}=x_i^j\), often called a Vandermonde matrix; its
invertibility follows from the interpolation argument just given.

\begin{proposition}[Vanishing on a domain]\label{prop:vanishing}
A polynomial \(f\) is zero at every point of \(H\) exactly when
\(Z_H\mid f\). Moreover,
\[
 L_i(X)=\frac{Z_H(X)}{(X-x_i)Z_H'(x_i)}.
\]
The quotient by \(X-x_i\) in this expression is a polynomial.
\end{proposition}
\begin{proof}
Remove the factor \(X-x_0\). At any other root, this factor is
nonzero, so the remaining polynomial still vanishes there. Repeat
to remove all factors. The reverse implication follows by evaluation.
The product rule gives
\(Z_H'(x_i)=\prod_{j\ne i}(x_i-x_j)\ne0\), proving the formula.
\end{proof}

Polynomial division now has an evaluation interpretation. The unique
remainder of \(f\) modulo \(Z_H\) is its degree-less-than-\(n\)
interpolant on \(H\), because subtraction of a multiple of \(Z_H\)
does not change those values. This is also the precise source of
aliasing: values on \(H\) determine a remainder, not an unrestricted
high-degree polynomial.

\subsection{Random polynomial masks and threshold observations}

Choose the \(d\) coefficients of a polynomial of degree less than
\(d\) independently and uniformly in \(\F_q\). At any \(k\le d\)
distinct points its evaluations are jointly uniform in \(\F_q^k\).
Indeed, restricting the evaluation map to the first \(k\) coefficients
already gives an invertible interpolation matrix, so the full map is
surjective. Apply Theorem~\ref{thm:affine-mask}.

Suppose instead that a polynomial must vanish on a fixed set \(H\).
Write it as \(Z_H(X)r(X)\), where \(r\) has \(d\) uniform
coefficients. At \(k\le d\) distinct points outside \(H\), its
evaluations remain jointly uniform: multiplication by the nonzero
values \(Z_H(x_i)\) is an invertible diagonal map. At a point inside
\(H\), the output is necessarily zero. Repeated observations at the
same point are not additional independent masks; their evaluation
rows coincide. Linear observations already made must be accounted for
using Lemma~\ref{lem:conditional-mask}.

\begin{lemma}[A hidden constant coefficient]\label{lem:threshold-polynomial}
Fix a value \(s\in\F_q\) and let
\(f(X)=s+\sum_{i=1}^{t-1}a_iX^i\), with independent uniform
coefficients. For any \(k<t\) distinct nonzero evaluation points,
the \(k\) evaluations have a uniform joint distribution independent
of \(s\). Conversely, evaluations at \(t\) distinct points
determine the complete polynomial and hence \(s\).
\end{lemma}
\begin{proof}
For the first claim, the columns corresponding to
\(X,\ldots,X^k\) give a matrix equal to the invertible
Vandermonde matrix for \(1,\ldots,X^{k-1}\), with each row
multiplied by its nonzero evaluation point. It has full row rank.
The affine-mask theorem removes dependence on the constant shift.
The second claim is interpolation.
\end{proof}

When reconstruction needs only the constant coefficient, interpolation
gives the explicit formula
\[
 f(0)=\sum_i \lambda_i f(x_i),\qquad
 \lambda_i=\prod_{j\ne i}\frac{-x_j}{x_i-x_j}.
\]
The index set in these products is the exact set of observations being
used. Changing that set changes the coefficients. The nonzero-point
condition in the hiding lemma is essential: an observation at zero
reveals \(s\) immediately.

\subsection{Several variables and random combinations}

A polynomial in \(m\) variables is a finite sum of monomials
\(cX_1^{e_1}\cdots X_m^{e_m}\). Its total degree is the largest
sum \(e_1+\cdots+e_m\) with nonzero coefficient. We may regard it
as a polynomial in the last variable whose coefficients are
polynomials in the others.

\begin{theorem}[Schwartz--Zippel]\label{thm:schwartz-zippel}
For a fixed nonzero polynomial of total degree at most \(d\), and
independent uniform \(x_1,\ldots,x_m\) in a finite nonempty set
\(S\subseteq\F\), the probability of a zero evaluation is at most
\(\min(1,d/|S|)\).
\end{theorem}
\begin{proof}
The case \(m=1\) is the root bound. For the induction step, let
\(k\) be the degree in \(X_m\) and \(g\ne0\) its leading
coefficient polynomial. Its total degree is at most \(d-k\), so
the probability that it vanishes on the first \(m-1\) choices is
at most \((d-k)/|S|\). Otherwise the remaining univariate
polynomial is nonzero of degree \(k\), and its conditional root
probability is at most \(k/|S|\), because \(x_m\) is still
independent uniform. Add the two bounds and cap at one.
\end{proof}

To combine fixed scalar equations \(e_0=0,\ldots,e_{k-1}=0\),
choose a fresh scalar \(\eta\) and test
\(\sum_i\eta^i e_i=0\). If at least one equation is false, this
is a nonzero polynomial of degree at most \(k-1\) in \(\eta\).
Independent uniform coefficients would instead give error exactly
\(1/q\) over \(\F_q\), by solving for the coefficient of a
nonzero \(e_i\). The two methods have different bounds.

For polynomial equations, first apply the same reasoning to a
coefficient that is nonzero in at least one equation; then account
for the later random evaluation of the resulting polynomial. These
are separate bad events. Random aggregation does not turn a single
root-bound calculation into a security proof for the complete
interactive protocol in which the equations arose.

\section{Evaluation domains and fast transforms}\label{sec:rootsofunity}

Interpolation works on any distinct points, but a geometric sequence
of points makes both evaluation and interpolation much faster. We
construct those sequences and the transforms rather than treating
the fast Fourier transform, abbreviated FFT, as an unexplained operation.

\subsection{Roots of unity and their existence}

An element \(\omega\in\F^\times\) is a primitive \(n\)-th root
of unity if its multiplicative order is exactly \(n\). Its powers
\(1,\omega,\ldots,\omega^{n-1}\) are then distinct, form a
subgroup, and are roots of \(X^n-1\). The root bound proves that
there are no further roots and hence
\(Z_H(X)=X^n-1\) for this set \(H\).

\begin{theorem}[Cyclicity of a finite multiplicative subgroup]
\label{thm:field-cyclic}
Every finite subgroup of a field's multiplicative group is cyclic.
In particular, \(\F_q^\times\) has an element of order \(q-1\).
\end{theorem}
\begin{proof}
Let \(M\) be the least common multiple of the orders of all elements.
Factor \(M\) into its maximal prime-power factors. For each such
factor \(\ell^a\), some element has
order divisible by \(\ell^a\). Raising that element to the part of
its order not divisible by \(\ell\) gives an element of order
\(\ell^a\). This follows directly from scalar reduction: for an
element of order \(r\), the order of its \(k\)-th power is
\(r/\gcd(r,k)\).

Two integers are coprime when their greatest common divisor is one.
The product of commuting elements of coprime orders has order the
product. If \((ab)^t=1\), then \(a^t=b^{-t}\) lies in the
intersection of the two cyclic subgroups. Its order divides both
coprime orders by Lagrange, so it is the identity; each order therefore
divides \(t\). Multiplying the elements constructed for all prime
powers yields an element of order \(M\).

Every member of the original subgroup is a root of \(X^M-1\), so
its size is at most \(M\) by the root bound. The element just
constructed already contributes \(M\) distinct powers. Thus the
subgroup consists exactly of those powers.
\end{proof}

There is a primitive \(n\)-th root in \(\F_q\) exactly when
\(n\mid(q-1)\): necessity follows from Lagrange, and for
sufficiency raise a multiplicative generator to \((q-1)/n\).
The largest \(s\) for which \(2^s\mid(q-1)\) is the
two-adicity of the field. An advertised root of order \(2^s\) can
be checked without finding a full generator: verify
\(\omega^{2^s}=1\) and \(\omega^{2^{s-1}}\ne1\) when \(s\ge1\).
When \(s=0\), the unique element of order one is simply 1. For general
\(n\), verify \(\omega^n=1\) and
\(\omega^{n/\ell}\ne1\) for each prime divisor \(\ell\) of
\(n\). Every proper divisor of \(n\) divides some \(n/\ell\),
which proves the test.

\subsection{The transform and its inverse}

The characteristic of a finite field is the least positive integer
\(\ell\) for which \(\ell\cdot1=0\). It exists because the
additive group is finite. It is prime: a factorisation
\(\ell=ab\) with \(1<a,b<\ell\) would produce two nonzero
elements \(a\cdot1,b\cdot1\) with zero product. Its additive
subgroup generated by one has order \(\ell\), so Lagrange shows
\(\ell\mid q\) for a field of size \(q\).

Let \(n\mid(q-1)\) and let \(\omega\) have order \(n\).
The field element represented by \(n\) is nonzero and invertible:
the characteristic prime divides \(q\), not \(q-1\). For an
integer \(k\), a geometric sum gives
\begin{equation}\label{eq:orthogonality}
 \sum_{j=0}^{n-1}\omega^{jk}=
 \begin{cases}n,&n\mid k,\\0,&n\nmid k.\end{cases}
\end{equation}
In the second case multiply the sum by \(1-\omega^k\ne0\);
the product is \(1-\omega^{nk}=0\). This proves the identity.

For coefficients \(a_0,\ldots,a_{n-1}\), the transform is
\(y_j=\sum_i a_i\omega^{ij}\), evaluation at \(\omega^j\).
Its inverse is
\begin{equation}\label{eq:inverse-transform}
 a_i=n^{-1}\sum_{j=0}^{n-1}y_j\omega^{-ij}.
\end{equation}
Substitute the expression for \(y_j\), interchange the two sums,
and use \eqref{eq:orthogonality}. Only the term with the matching
coefficient index remains. In particular, the Lagrange basis here is
\[
 L_j(X)=\frac1n\sum_{i=0}^{n-1}\omega^{-ij}X^i
       =\frac{\omega^j}{n}\frac{X^n-1}{X-\omega^j}.
\]
The second expression follows from the derivative
\((X^n-1)'=nX^{n-1}\) and Proposition~\ref{prop:vanishing}.
At \(X=\omega^j\) it denotes the polynomial quotient, not a
numerical division by zero.

\subsection{The radix-two algorithm}

For even \(n\), separate the even and odd coefficients:
\[
 f(X)=f_{\rm even}(X^2)+Xf_{\rm odd}(X^2).
\]
The root \(\omega^2\) has order \(n/2\). Compute the two
length-\(n/2\) transforms to obtain values \(E_j,O_j\) at
\(\omega^{2j}\). For \(0\le j<n/2\), combine them as
\[
 y_j=E_j+\omega^jO_j,\qquad
 y_{j+n/2}=E_j-\omega^jO_j.
\]
The identity \(\omega^{n/2}=-1\) follows because it has order two,
and the only roots of \(X^2-1\) in a field of odd characteristic
are \(1,-1\). These formulas compute the two required evaluations.
When \(n=2^k\), repeat the split down to length one, where a
constant evaluates to itself.

At each of \(k\) levels the total number of field operations is
bounded by a fixed multiple of \(n\); the result costs
\(O(n\log n)\), instead of evaluating one degree-\(n-1\)
polynomial at \(n\) points by separate quadratic total work.
The inverse uses the
same algorithm with root \(\omega^{-1}\), followed by multiplication
of each result by \(n^{-1}\). An in-place implementation may
reorder entries; the output contract remains the explicit indexing
in \eqref{eq:inverse-transform}.

\begin{example}[A complete transform]\label{ex:fft}
In \(\F_{97}\), the element 22 has order four, with powers
\((1,22,96,75)\). The polynomial
\(3+5X+7X^2+11X^3\) evaluates on this domain to
\((26,58,91,31)\). Applying the inverse transform recovers
\((3,5,7,11)\). Both directions are checked independently in the
accompanying script.
\end{example}

\subsection{Products, rotations and cosets}

To multiply polynomials of degrees at most \(a,b\), choose an
available transform length \(N>a+b\), pad both coefficient lists
with zeros, evaluate them, multiply corresponding values, and
interpolate. The true product has degree less than \(N\), so
interpolation uniquely recovers it. If \(N\le a+b\), the same
procedure instead returns the remainder modulo \(X^N-1\).
Calling it the full product would be incorrect.

For a polynomial \(f\), the rotation \(f(\omega X)\) evaluates
on \(H\) to the cyclic shift \(f(\omega^{j+1})\). Its
coefficient of \(X^i\) is \(a_i\omega^i\). The wrap at the
last row is an algebraic property of the domain, not a decision about
which circuit rows should participate in a constraint.

A multiplicative coset is \(gH=\{g\omega^j\}\), with
\(g\ne0\). Evaluating \(f\) on it is the ordinary transform
of coefficients \(a_ig^i\). Interpolation reverses this scaling.
Its vanishing polynomial is \(X^n-g^n\). If \(g\notin H\),
then \(g^n\ne1\), and the coset is disjoint from \(H\).

Suppose a quotient \(t=f/(X^n-1)\) exists and has degree less
than \(N\), where \(n\mid N\) and a size-\(N\) subgroup
\(H_N\) is available. Choose \(g\notin H_N\). Every point of
\(gH_N\) lies outside \(H\), because otherwise \(g\) would
belong to \(H_N\). Thus
\[
 t(x)=\frac{f(x)}{x^n-1}\qquad(x\in gH_N)
\]
has no zero denominator and can be interpolated on that coset.
The bound needed for this interpolation is on \(\deg t\), not
on \(\deg f\). To obtain the numerator values one may evaluate
its constituent low-degree polynomials on the coset and perform the
specified pointwise expression; one need not interpolate the
high-degree numerator itself in a space that is too small for it.

There is a converse trap. Pointwise division on a coset followed by
interpolation always returns some polynomial, even when \(f\) is
not divisible by \(X^n-1\). That does not prove the polynomial
identity \(f=(X^n-1)t\). Divisibility must be established separately
or tested with the appropriate degree and challenge conditions.

\subsection{Chunks and recombination}

Every polynomial of degree less than \(kn\) decomposes uniquely as
\[
 t(X)=\sum_{j=0}^{k-1}X^{jn}t_j(X),\qquad \deg t_j<n,
\]
by grouping consecutive blocks of \(n\) coefficients. At a point
\(z\), its value is \(\sum_j z^{jn}t_j(z)\). This is a
deterministic identity, with no randomness or probability loss.
It allows a large quotient to be represented by several polynomials
with a common smaller degree bound. A later protocol may commit to
the chunks separately; such a commitment does not alter the
recombination formula or the need for the degree bounds.

\section{Pasta curves and their encodings}\label{sec:ellipticcurves}

We now construct the point arithmetic used throughout the cryptographic
volumes. Two distinctions remain visible: coordinates are not scalars,
and an algebraic formula valid away from exceptional inputs is not a
total group operation. The finite-field algorithms already developed
make every ordinary arithmetic step explicit.

\subsection{The curve and the complete affine operation}

An affine coordinate pair means two ordinary field coordinates
\((x,y)\), without identifying different pairs by a scaling
rule. We first specify the curve and its operation in those
coordinates, treating its extra identity as a separate object.

For a field \(\F_p\) with prime \(p>3\), choose \(A,B\in\F_p\)
with \(4A^3+27B^2\ne0\). Consider the finite coordinate pairs
satisfying
\begin{equation}\label{eq:weierstrass}
 y^2=x^3+Ax+B,
\end{equation}
and add a separate point \(\mathcal O\), called the point at
infinity. This equation is called short-Weierstrass form. The
nonvanishing condition is a hypothesis of the group theorem
below. To see its algebraic meaning, a repeated root \(u\)
of \(X^3+AX+B\) would satisfy
\(3u^2+A=0\) and \(u^3+Au+B=0\), hence
\(A=-3u^2,B=2u^3\), giving \(4A^3+27B^2=0\).

Negate a finite point by \(-(x,y)=(x,-y)\), and set
\(-\mathcal O=\mathcal O\). Define addition with the identity
by \(P+\mathcal O=\mathcal O+P=P\). For finite points
\(P=(x_1,y_1)\), \(Q=(x_2,y_2)\), use the following cases:
\begin{enumerate}
\item If \(x_1=x_2\) and \(y_1=-y_2\), return \(\mathcal O\).
      This includes doubling a point with \(y_1=0\).
\item Otherwise, if \(P=Q\), set
      \(\lambda=(3x_1^2+A)/(2y_1)\).
\item Otherwise, set \(\lambda=(y_2-y_1)/(x_2-x_1)\).
\end{enumerate}
In the last two cases return
\begin{equation}\label{eq:curve-add}
 x_3=\lambda^2-x_1-x_2,\qquad
 y_3=\lambda(x_1-x_3)-y_1.
\end{equation}
The cases cover all inputs. With equal \(x\)-coordinates, the curve
equation gives \(y_1^2=y_2^2\); factoring their difference shows
the points are equal or opposite. The excluded cases ensure every
displayed denominator is nonzero.

These formulas come from intersecting a line with the cubic. Substitute
the line \(y=\lambda(x-x_1)+y_1\) into
\eqref{eq:weierstrass}. After moving all terms to one side, the
coefficient of \(x^2\) is \(-\lambda^2\). Two roots are
\(x_1,x_2\), so the third is \(\lambda^2-x_1-x_2\), by
expanding the product of the three linear factors. Reflection in the
\(x\)-axis gives \(y_3\). For equal inputs, the slope is chosen
so that \(x_1\) is a repeated root of the substituted polynomial;
the formal derivative condition gives \(2y_1\lambda=3x_1^2+A\).
This proves that the returned finite point is on the curve and
explains both ordinary addition and doubling without division by a
vanishing denominator.

\begin{theorem}[Imported elliptic-curve group theorem]
\label{thm:curve-group-boundary}
For \(p>3\) and \(4A^3+27B^2\ne0\), the points of
\eqref{eq:weierstrass}, together with \(\mathcal O\), form an
abelian group under the complete operation just defined.
\end{theorem}

Closure was established above; the identity, inverses and symmetry of
the construction are explicit. Associativity, including all exceptional
cases, is the imported part of this theorem. The line-intersection
picture does not prove it, and checking a few triples numerically does
not replace it. We use the standard short-Weierstrass group theorem
also assumed by the protocol specification's section ``Pallas and
Vesta''. This is an explicit mathematical boundary of this volume,
not a computational hardness assumption and not a claim that the
preceding coefficient calculation proves the entire theorem.

\subsection{Incomplete addition and circuit equations}

The incomplete operation \(P\boxplus Q\) is defined only for finite
points with different \(x\)-coordinates. It uses the last slope
formula and \eqref{eq:curve-add}; on any other input it returns
\(\bot\). Thus it agrees with complete addition wherever defined,
but it rejects \(P=Q\), \(P=-Q\), and any identity input.
Equality of two finite \(x\)-coordinates is equivalent to equality
up to negation, as proved above.

One can replace the division in its slope by the equation
\begin{equation}\label{eq:incomplete-slope}
 (x_2-x_1)\lambda=y_2-y_1.
\end{equation}
With \(x_2\ne x_1\), this equation uniquely determines the
correct slope. With \(P=Q\), it becomes \(0=0\) and leaves
\(\lambda\) unconstrained. Consequently, polynomial equations
obtained by clearing denominators do not automatically enforce the
domain of an incomplete operation. A circuit may need a nonzero
check, or its surrounding security argument may need to account
explicitly for exceptional witnesses allowed by the implemented
relation. A description of a partial function and a description of
its circuit constraints must not silently be identified.

A nonzero check can introduce \(z\) and enforce \(xz=1\): an
inverse exists precisely when \(x\ne0\). A conditional nonzero
check with a Boolean flag \(b\), meaning an element constrained
to zero or one, instead needs a correctly specified
conditional relation, since \(xz=b\) alone does not force \(x=0\)
when \(b=0\). Such details are statements about equations, not
minor implementation conventions.

\subsection{Projective representation and scalar evaluation}

Affine formulas require inversions. A projective coordinate system
introduces a shared denominator and identifies scaled
representatives. In ordinary projective coordinates, a point
is an equivalence class of nonzero triples \((X,Y,Z)\), where
\((X,Y,Z)\) and \((uX,uY,uZ)\) represent the same point for
every nonzero \(u\). For \(Z\ne0\), the corresponding affine
coordinates are \((X/Z,Y/Z)\). The homogeneous curve equation,
meaning that every term has the same total degree, is
\[
 Y^2Z=X^3+AXZ^2+BZ^3.
\]
It is unchanged by the scaling because every term acquires a
factor \(u^3\). With \(Z=0\), the equation forces \(X=0\),
and the nonzero-triple condition leaves the single equivalence
class of \((0,1,0)\). This constructs the extra point at
infinity in the same coordinate language as the finite points.

Jacobian coordinates use a different scaling rule, chosen to
reduce arithmetic costs. A finite point is represented by
\((X,Y,Z)\), with \(Z\ne0\), meaning
\(x=X/Z^2\), \(y=Y/Z^3\). Its equation becomes
\[
 Y^2=X^3+AXZ^4+BZ^6.
\]
Scaling the triple to \((u^2X,u^3Y,uZ)\), with \(u\ne0\), does
not change the represented point. This is a weighted scaling:
the powers of \(u\) are two, three and one rather than all
one. For the algorithms here the identity is a separate tagged
value. An implementation may instead reserve particular
zero-\(Z\) triples for it; that internal convention is not
an ordinary affine point and must be handled before division.

To see how denominator removal works, consider two finite triples.
Set
\[
 U_1=X_1Z_2^2,\quad U_2=X_2Z_1^2,\qquad
 S_1=Y_1Z_2^3,\quad S_2=Y_2Z_1^3,
\]
and \(H=U_2-U_1\), \(R=S_2-S_1\). When \(H\ne0\), return
\begin{align*}
 X_3&=R^2-H^3-2U_1H^2,\\
 Y_3&=R(U_1H^2-X_3)-S_1H^3,\\
 Z_3&=HZ_1Z_2.
\end{align*}
Dividing these expressions by \(Z_3^2,Z_3^3\) respectively
recovers \eqref{eq:curve-add}, because
\(x_2-x_1=H/(Z_1^2Z_2^2)\) and
\(y_2-y_1=R/(Z_1^3Z_2^3)\). This is an explicit
inversion-free ordinary-addition formula, not a complete formula
for \(H=0\). Doubling, opposites and identities require their own
cases or a separately verified complete projective formula.

For correctness, one may always convert to affine and use the complete
operation already specified. For performance, an implementation
retains projective coordinates through a sequence of additions and
normalises once at the end. Batch normalisation can use the simultaneous
inversion algorithm in \S\ref{sec:intmod}. No choice of coordinates
alters the distinction between the coordinate field and the scalar
modulus.

\subsection{Quadratic residues and square-root extraction}

A nonzero element \(a\in\F_p\) is a square if \(a=y^2\) for
some \(y\). A nonzero square has exactly two square roots \(y,-y\),
because \(X^2-a\) has at most two roots and the two are distinct
when \(p\) is odd. The squaring map on \(\F_p^\times\) is
therefore two-to-one onto a set of size \((p-1)/2\). By cyclicity,
writing \(a=g^k\) for a generator \(g\), the squares are exactly
the even powers. The element \(g^{(p-1)/2}\) is \(-1\), so
\begin{equation}\label{eq:euler-criterion}
 a^{(p-1)/2}=\begin{cases}1,&a\text{ is a nonzero square},\\
                         -1,&a\text{ is a nonsquare}.
             \end{cases}
\end{equation}
This is Euler's criterion, including an efficient test by exponentiation.
The zero input has the unique square root zero and is handled separately.

Here is a constructive algorithm valid for every odd prime. Write
\(p-1=2^s d\), with \(d\) odd, and choose a nonsquare \(z\),
verified with \eqref{eq:euler-criterion}. For a nonzero input \(a\),
first reject unless Euler's test gives one. Initialise
\[
 x=a^{(d+1)/2},\qquad t=a^d,\qquad c=z^d,\qquad m=s.
\]
While \(t\ne1\), find the least \(i\ge1\) with
\(t^{2^i}=1\), set \(b=c^{2^{m-i-1}}\), and update
\[
 x\leftarrow xb,\quad t\leftarrow tb^2,\quad
 c\leftarrow b^2,\quad m\leftarrow i.
\]
Return \(x\). This is the Tonelli--Shanks algorithm.

\begin{proposition}[Square-root algorithm correctness]
\label{prop:tonelli-shanks}
For the accepted inputs above, the loop is defined, terminates and
returns a square root of \(a\). It uses
\(O(\log p+s^2)\) field multiplications after a verified nonsquare
has been supplied. A search for that nonsquare is additional work;
independent uniform trials need two trials on average, with an
\(O(\log p)\) Euler test per trial. This is an expected bound,
not a fixed worst-case bound for deterministic sequential search.
\end{proposition}
\begin{proof}
Maintain \(x^2=at\), \(\ord(c)=2^m\), and
\(\ord(t)\mid2^{m-1}\). Initially the first identity is a direct
power calculation. Euler's criterion gives
\(c^{2^{s-1}}=-1\), so \(c\) has order \(2^s\), and
\(t^{2^{s-1}}=1\) since \(a\) is a square. Thus all
invariants hold.

When \(t\ne1\), its order is \(2^i\) with \(1\le i<m\), so
the exponent defining \(b\) is nonnegative. The element \(b\)
has order \(2^{i+1}\), and \(b^2\) has order \(2^i\).
Both \(t\) and \(b^2\), raised to \(2^{i-1}\), equal the
unique order-two element \(-1\). Their product consequently has
order dividing \(2^{i-1}\). The updated variables satisfy all
three invariants with the strictly smaller value \(m=i\).
Termination follows, and the first invariant gives \(x^2=a\)
when \(t=1\). Initial exponentiations cost \(O(\log p)\);
each of at most \(s\) iterations uses at most \(O(s)\)
squarings or multiplications.
\end{proof}

A nonsquare can be found by testing candidates, or by independent
uniform sampling from \(\F_p^\times\), which succeeds with
probability one half per trial. The variable-time algorithm above
specifies an arithmetic result; an implementation handling secrets
must separately prevent execution time or memory accesses from
revealing sensitive choices. A final check \(x^2=a\) is a useful
implementation assertion but does not replace input validation.

\subsection{Concrete parameters and the two field roles}

Define
\begin{align*}
 p_{\Pallas}&=2^{254}+45560315531419706090280762371685220353,\\
 p_{\Vesta}&=2^{254}+45560315531506369815346746415080538113.
\end{align*}
The curve \(\Ep\) has equation \(y^2=x^3+5\) over
\(\F_{p_{\Pallas}}\); the curve \(\Ev\) has the same equation
over \(\F_{p_{\Vesta}}\). The short-Weierstrass nonvanishing
condition is satisfied because \(27\cdot25\ne0\) in either field.

\begin{assumption}[Imported, verifiable parameter facts]
\label{ass:pasta-parameters}
Both displayed integers are prime, and the complete point groups have
orders
\[
 |\Ep(\F_{p_{\Pallas}})|=p_{\Vesta},\qquad
 |\Ev(\F_{p_{\Vesta}})|=p_{\Pallas}.
\]
These are mathematical parameter facts imported from the Zcash
protocol specification, ``Pallas and Vesta'', not assumptions of
computational intractability.
\end{assumption}

The specification is available at
\href{https://zips.z.cash/protocol/protocol.pdf}{the primary protocol
source}. The inspected implementation is \texttt{pasta\_curves}
version 0.5.2, with moduli in \texttt{src/fields/fp.rs} and
\texttt{src/fields/fq.rs}, and group/type definitions in
\texttt{src/curves.rs}. These identify the rule set being explained;
they do not by themselves establish the activation of any network
upgrade.

Under these parameter facts, every nonidentity point generates its
group and both cofactors are one. The Pallas base field, used for
coordinates, is \(\F_{p_{\Pallas}}\), but its scalar field is
\(\F_{p_{\Vesta}}\). Vesta reverses these roles. This is the
Pasta cycle: the base field of each curve is the scalar field of the
other. The cycle makes Pallas coordinate operations expressible
directly in arithmetic over Vesta's scalar field. It does not identify
Pallas scalars with Pallas coordinates, and it does not make all
operations in an algorithm using both curves become operations
in a single field.

Both fields have two-adicity exactly 32: divide the corresponding
\(p-1\) by \(2^{32}\), and the remaining integer is odd. The
checker verifies this, as well as Euler's criterion showing that
\(-13\) and 5 are nonsquares in both fields. The point \((-1,2)\)
satisfies \(4=(-1)^3+5\), so it is a nonidentity point in each
group. The checker also verifies that multiplication by the
published opposite-field order returns the identity. That check
confirms a scalar-multiplication identity; it is not a primality
certificate or a proof that there are no additional curve points.
The imported order facts remain explicit.

\subsection{Compressed points and coordinate extraction}

Let \(p\) be the relevant base modulus. For a finite point
\(P=(x,y)\), write \(\widetilde y\) for the low bit of the
canonical integer representative of \(y\). The star encoding is
\begin{equation}\label{eq:star-encoding}
 P^\star=\LE_{256}(\val(x)+2^{255}\widetilde y),\qquad
 \mathcal O^\star=\LE_{256}(0).
\end{equation}
Here \(\val(x)\) denotes the canonical integer in \([0,p)\).
Packing each consecutive eight bits into a byte gives the 32-byte
representation. The top bit carries the sign selector; the lower
255 bits carry the \(x\)-coordinate.

Decoding performs the following exact checks:
\begin{enumerate}
\item Require exactly 32 bytes. Extract the top bit \(s\), and
      interpret the remaining 255 bits as an integer \(u\).
\item Reject if \(u\ge p\). Otherwise set \(x=[u]_p\).
\item If \(u=0\) and \(s=0\), return \(\mathcal O\).
\item Compute a square root \(y\) of \(x^3+5\). Reject if no
      root exists. Select \(y\) or \(-y\) so that the canonical
      integer's low bit equals \(s\), then return \((x,y)\).
\end{enumerate}
This is the decoding behaviour in \texttt{pasta\_curves} 0.5.2,
\texttt{src/curves.rs}, implementing the specification's ``Pallas and
Vesta'' encoding. In particular, reduction of a noncanonical
\(u\ge p\) would be incorrect.

There are no finite points with \(y=0\), because such a point would
be a nonidentity element of order two in a group of odd prime order.
The two square roots in step four are therefore distinct and have
opposite parity, since their canonical representatives sum to the odd
integer \(p\). There are also no finite points with \(x=0\),
because 5 is a nonsquare in these fields. Consequently, the identity
encoding does not collide with a finite point, and the input with
only the top bit set is invalid rather than another identity encoding.
Every valid group point has exactly one valid byte representation.

Define coordinate extraction by
\(\operatorname{Extract}(\mathcal O)=0\) and
\(\operatorname{Extract}(x,y)=\val(x)\). Extending it to a partial
operation preserves failure: \(\operatorname{Extract}(\bot)=\bot\),
not zero. Equality of extracted coordinates for two finite points
means the points are equal or opposite; it does not generally mean
the full points are equal. For Pasta, extracted zero identifies the
identity, because finite \(x=0\) points do not exist. These exact
facts are required whenever a commitment is shortened to one
coordinate.

\subsection{A small curve and the limits of examples}

Over \(\F_{19}\), the equation \(y^2=x^3+5\), together with
the identity, has 27 points. Direct enumeration of all 361 coordinate
pairs establishes this count in the checker. The point \((0,9)\)
doubles to \((0,10)\), its negative, and its third multiple is the
identity. This illustrates the complete operation and a non-prime
curve order. It does not use Pasta's compressed-identity convention,
whose disjointness relied on the absence of finite \(x=0\) points.

No small example establishes a security level. Even the exact Pasta
group-order facts would not prove that discrete logarithms are hard.
They establish the algebraic domains on which a later security
assumption can be stated. The imported group theorem and parameter
facts are the only substantial unproved mathematical inputs to this
curve construction; all encoding operations and ordinary arithmetic
algorithms used here have been given explicitly.

\section{Byte fields for AES}\label{sec:byte-fields}

A byte-oriented algorithm may use a field that is not a prime field.
For AES, bytes represent polynomials with binary coefficients, not
integers modulo 256. This final construction supplies exactly that
field and its invertible operations.

\subsection{Binary polynomial arithmetic}

Use \(\F_2=\{0,1\}\), where addition and subtraction coincide
with exclusive-or. Represent a byte \(a\) by
\[
 a(X)=a_0+a_1X+\cdots+a_7X^7,
\]
where \(a_i\) are the bits of its integer encoding. Define
\(M(X)=X^8+X^4+X^3+X+1\). Add polynomials coefficientwise and
multiply them by convolution, then reduce the result modulo \(M\)
using polynomial division. The result again has degree less than
eight and so fits in one byte. Replacing either input by a congruent
polynomial does not change the result, because multiples of \(M\)
remain multiples under addition and multiplication. Associativity
and distributivity therefore descend from \(\F_2[X]\).

A polynomial of positive degree is irreducible if it is not a
product of two polynomials of smaller positive degree. A reducible
degree-eight polynomial has a factor of degree at most four: in
any factorisation, at least one of the two degrees is at most four.
Since the only nonzero leading coefficient in \(\F_2\) is one,
all candidate factors are monic, meaning that their leading
coefficient is one. There are
\(2+4+8+16=30\) monic polynomials of degrees one through four.
Dividing \(M\) by each is a complete finite irreducibility test.
The accompanying checker performs all 30 divisions and verifies
that none has zero remainder. The hexadecimal representation of
the full modulus is \(\mathtt{0x11b}\); the ninth bit records
the \(X^8\) term and is not itself an eight-bit field element.

\begin{proposition}[The AES field]\label{prop:aes-field}
The 256 degree-less-than-eight binary polynomials, with the
operations above, form a field.
\end{proposition}
\begin{proof}
We have established the arithmetic laws inherited from polynomials.
To invert a nonzero polynomial \(a\) of degree less than eight,
run Euclid's algorithm on \(a,M\), using polynomial division.
The last nonzero remainder is a common divisor and a polynomial
combination \(ua+vM\), by exactly the coefficient bookkeeping
in the integer algorithm. A nonconstant common divisor would
factor the irreducible polynomial \(M\); its degree cannot be
eight because it also divides \(a\). Thus the greatest common
divisor is the constant one. Reducing the identity
\(ua+vM=1\) modulo \(M\) gives the inverse \(u\bmod M\).
\end{proof}

The field contains a copy of \(\F_2\): send zero and one
to their constant-polynomial representatives. This map
preserves one, addition and multiplication, and is injective.
Its image is therefore a subfield, a subset that is itself a
field under the inherited operations. We may identify that
image with the original \(\F_2\).

A field containing a specified copy of another field is a
field extension of that smaller field. Thus the byte field is
an extension of \(\F_2\), not merely a larger set carrying
the same names for some elements. Multiplication by a
constant polynomial makes it a vector space over \(\F_2\).
Its eight representatives \(1,X,\ldots,X^7\) form a basis:
every reduced polynomial is their linear combination, and
a nonzero combination of degree less than eight cannot be a
multiple of the degree-eight modulus. The dimension of this
vector space is the degree of the field extension, here eight.
Consequently it has exactly \(2^8\) elements, one for each
choice of eight binary coefficients. This derives both the
extension terminology and the size from the constructed
operations; no unspecified extension arithmetic is being used.

We denote this field by \(\F_{2^8}\). Its concrete byte
representation depends on the chosen modulus and basis.
A field isomorphism preserves its mathematical operations,
as defined in \S\ref{sec:intmod}, but need not preserve the
integer labels of the bytes. Consequently another byte basis
can give a different encoded multiplication table. The symbol
\(\mathtt{0x02}\) represents \(X\), not the prime-field
integer two, which is zero in characteristic two.

\subsection{Constructive multiplication and inversion}

Multiplication by \(X\) shifts the byte left by one bit. If the
old highest bit was one, the shift introduced \(X^8\), which
reduces to \(X^4+X^3+X+1\). Therefore the byte operation is
\[
 \operatorname{xtime}(a)=
 ((a\ll1)\bmod256)\oplus
 \begin{cases}\mathtt{0x1b},&a\ge128,\\0,&a<128.\end{cases}
\]
To multiply \(a\) by \(b\), initialise a zero accumulator and
scan the eight bits of \(b\) from low to high. Exclusive-or the
current value of \(a\) into the accumulator for each set bit,
and apply \(\operatorname{xtime}\) to \(a\) after each step.
This is polynomial convolution followed by reduction, so it proves
the operation implements the field product.

The nonzero multiplicative group has order 255. Lagrange gives
\(a^{255}=1\), so \(a^{-1}=a^{254}\) for nonzero \(a\).
Exponentiation using the byte multiplication just specified is a
second complete inversion algorithm. A later substitution operation
may define an extended inverse that maps zero to zero. That is a
separate total function, not a claim that zero has a multiplicative
inverse.

\begin{example}[Byte arithmetic]\label{ex:aes-byte}
The field product \(\mathtt{0x57}\cdot\mathtt{0x83}\) is
\(\mathtt{0xc1}\), and the inverse of \(\mathtt{0x53}\) is
\(\mathtt{0xca}\). The checker evaluates both through the explicit
polynomial arithmetic and checks \(aa^{254}=1\) for all 255
nonzero byte values.
\end{example}

An affine map on the eight bits means a binary linear transformation
followed by exclusive-or with a fixed byte. To test its invertibility,
form the eight-by-eight matrix of its effect on the coordinate bits
and apply elimination over \(\F_2\). A byte substitution formed
by the extended inverse followed by an invertible affine map is a
permutation: the extended inverse exchanges each nonzero element
with its inverse and fixes zero; composition preserves bijectivity.
This establishes the algebra needed to construct AES's substitution
step, without identifying byte-field arithmetic with ordinary
integer arithmetic or assuming a cipher's security from its
invertibility.

\subsection*{Verification boundary}

The numerical examples in this volume are produced and checked by
\texttt{check\_mathcrypto.py} alongside the rewrite sources. The
checks include field inverses and reduction bias, rank-one and
rank-two mask distributions, both directions of the four-point
transform, a complete small-curve point count, published Pasta
scalar identities, and exhaustive byte-field inverses and modulus
irreducibility. They do not certify Pasta primality or point counts,
prove the imported elliptic-curve associativity theorem, establish
cryptographic hardness, or certify constant-time implementations.
Those limits are part of the mathematical contract, not omissions
to be filled by trusting a worked example.
\end{document}
